\documentclass[11pt]{article}
\usepackage[margin=24mm]{geometry}
\usepackage[T1]{fontenc}
\usepackage{amsmath,amssymb,amsthm,mathtools,mathrsfs}
\usepackage[hidelinks]{hyperref}
\allowdisplaybreaks
\setlength{\emergencystretch}{3em}
\title{Original theta closure and arithmetic observation: complete literature receivers}
\author{}
\date{20 September 2026}
\begin{document}\maketitle
This edition proves three exact uses of the literature in the original
Split-Zero cohomology programme: the original Gamma evaluation kernel and
its three-vector arithmetic observation constraints; closure of the full
theta source in its specified strong topology; and reconstruction of the
actual arithmetic Gram from a single coefficient sequence, retaining the
finite Jacobi boundary correction. Each part contains its full derivation
and identifies the human source at the point of use. No sign is assigned
to an unevaluated projected current. These results do not close the RH
programme or assert existence of an off-line zero.

% BEGIN INDEXED LITERATURE EXACT RECEIVERS 20260920

\clearpage\begingroup
\part*{Exact Gamma kernel and the actual joint observed Gram}
\newtheorem{lgammatheorem}{Theorem}
\newtheorem{lgammalemma}[lgammatheorem]{Lemma}
\newtheorem{lgammaproposition}[lgammatheorem]{Proposition}
\newcommand{\lgammaC}{\mathbb C}
\newcommand{\lgammaR}{\mathbb R}
\newcommand{\lgammanorm}[1]{\left\lVert#1\right\rVert}
\newcommand{\lgammaim}{\operatorname{im}}
\newcommand{\lgammaremm}{\operatorname{rem}}



This is a standalone mathematical derivation.  All source masses,
coordinates, physical units, observation maps, original arithmetic boundary
vectors and affine relation spaces remain present.  The new finite Gamma
kernel formula is transported to the arithmetic source using the already
proved whole-polynomial comparison in CAI16--18 and OB13,18
\cite{lgamma:CAI,lgamma:OB}.  The receiving object is the actual three-vector Gram from
HG6--7 and HR2--5 \cite{lgamma:HG,lgamma:HR}, not a newly selected Gamma boundary triple.
The resulting constraints retain their positive-semidefinite remainders
and every complex cross term.  No projected current sign is assigned by
discarding an unevaluated remainder.

\section{The original Gamma measure and its exact norms}

Put
\begin{equation}\label{lgamma:eq:sigma}
 d\sigma(y)=\frac{|\Gamma(\tfrac14+iy/2)|^2}{2\pi}\,dy,
 \qquad M_\Gamma=\sqrt{2\pi}.
\end{equation}
The following calculation both fixes the mass and proves the polynomial
identities used below; no change of measure is made.
Euler's beta integral, with $a>0$, gives
\[
 |\Gamma(a+it)|^2
 =\Gamma(2a)\int_{\lgammaR}e^{itu}(2\cosh(u/2))^{-2a}\,du.
\]
Indeed this is the substitution $x=e^u$ in the beta integral on
$(0,\infty)$.  The function of $u$ is integrable.  The gamma product is
integrable as well, since Stirling's formula gives exponential decay
$e^{-\pi|t|}$ times a power.  Fourier inversion and evenness give
\[
 \int_{\lgammaR}|\Gamma(a+it)|^2e^{i\xi t}\,dt
 =2\pi\Gamma(2a)(2\cosh(\xi/2))^{-2a}.
\]
At $a=1/4$, substituting $y=2t$ in this integral, with the Jacobian
retained, proves
\begin{equation}\label{lgamma:eq:transform}
 \int_{\lgammaR}e^{i\xi y}\,d\sigma(y)
 =M_\Gamma(\cosh\xi)^{-1/2}.
\end{equation}
This substitution proves an integral identity and does not redefine the
measure or the polynomial coordinate.  In particular
$\sigma(\lgammaR)=M_\Gamma$.  The same gamma decay proves exponential
integrability for real $s$ with $|s|<\pi/2$.  Analytic continuation in
$|\Im\xi|<\pi/2$, with the branch positive at zero, therefore gives
\begin{equation}\label{lgamma:eq:laplace}
 \int_{\lgammaR}e^{sy}\,d\sigma(y)=M_\Gamma(\cos s)^{-1/2}.
\end{equation}

Define the monic polynomials in the original variable by
\begin{equation}\label{lgamma:eq:gen}
 F(y,t)=\sum_{n\ge0}p_n^\sigma(y)\frac{t^n}{n!}
 =(1+t^2)^{-1/4}e^{y\arctan t}.
\end{equation}
For sufficiently small real $t,u$, the cosine addition identity and
\eqref{lgamma:eq:laplace} give
\[
 \int_{\lgammaR}F(y,t)F(y,u)\,d\sigma(y)=M_\Gamma(1-tu)^{-1/2}.
\]
Exponential integrability justifies every derivative at $t=u=0$.
Comparing coefficients proves
\begin{equation}\label{lgamma:eq:norms}
 \int p_n^\sigma p_m^\sigma\,d\sigma
 =\delta_{nm}\gamma_n,
 \qquad \gamma_n=M_\Gamma n!(\tfrac12)_n.
\end{equation}
Also $(1+t^2)\partial_tF=(y-t/2)F$, so
\begin{equation}\label{lgamma:eq:monicrec}
 p_0^\sigma=1,\quad p_1^\sigma=y,\qquad
 p_{n+1}^\sigma=yp_n^\sigma-n(n-\tfrac12)p_{n-1}^\sigma.
\end{equation}
Thus the Gamma polynomial formulas are proved here directly.  Their
existing programme occurrence is CAI19 \cite{lgamma:CAI}; no result from CCM
is used in this derivation.

\section{Exact finite evaluation of the imaginary-point kernel}

\begin{lgammatheorem}\label{lgamma:thm:kernel}
For $N\ge0$, put $m_N=\lceil N/2\rceil$ and
\[
 \tau_m=\frac{(\tfrac12)_m(\tfrac34)_m}{m!(\tfrac14)_m}.
\]
For precisely the measure \eqref{lgamma:eq:sigma},
\begin{equation}\label{lgamma:eq:exactkernel}
 \boxed{K_N(i,i):=\sum_{n=0}^N\frac{|p_n^\sigma(i)|^2}{\gamma_n}
       =\frac{N+1}{M_\Gamma}\tau_{m_N}.}
\end{equation}
In particular, the exact reciprocal scalar is
\begin{equation}\label{lgamma:eq:beta}
 \boxed{\beta_N=K_N(i,i)^{-1}
 =\frac{\sqrt{2\pi}}{(N+1)\tau_{\lceil N/2\rceil}}.}
\end{equation}
The following explicit finite-degree bounds also hold:
\begin{align}
 K_N(i,i)&\le\frac{2N+1}{M_\Gamma}&& (N\ge0),\label{lgamma:eq:linear}\\
 K_N(i,i)&\le\frac{18N+9}{10M_\Gamma}&& (N\ge2).\label{lgamma:eq:betterlinear}
\end{align}
\end{lgammatheorem}
\begin{proof}
Write $p_n^\sigma(i)=i^na_n$.  Equation \eqref{lgamma:eq:monicrec} proves
\[
 a_0=a_1=1,\qquad
 a_{n+1}=a_n+n(n-\tfrac12)a_{n-1}>0.
\]
Inductively,
\begin{equation}\label{lgamma:eq:alternating}
 \frac{a_n}{a_{n-1}}=
 \begin{cases}n,&n\text{ odd},\\n-\tfrac12,&n\text{ even}.
 \end{cases}
\end{equation}
For the induction step, $a_n=na_{n-1}$ at odd $n$ gives
$a_{n+1}=(n+1/2)a_n$, while $a_n=(n-1/2)a_{n-1}$ at even $n$
gives $a_{n+1}=(n+1)a_n$.  The initial ratio is one.  Multiplication gives
\[
 a_{2m}=4^m(\tfrac12)_m(\tfrac34)_m,
 \qquad a_{2m+1}=(2m+1)a_{2m}.
\]
Let $u_n=a_n/\sqrt{\gamma_n}$.  Grouping the even and odd factors
in $(2m)!$ and $(1/2)_{2m}$ gives
\begin{equation}\label{lgamma:eq:terms}
 M_\Gamma u_{2m}^2=\tau_m,\qquad
 M_\Gamma u_{2m+1}^2=\frac{2m+1}{2m+\tfrac12}\tau_m,
 \qquad
 \frac{\tau_{m+1}}{\tau_m}
 =\frac{(m+\tfrac12)(m+\tfrac34)}{(m+1)(m+\tfrac14)}.
\end{equation}
The partial sums satisfy exactly
\[
 M_\Gamma K_{2m}=(2m+1)\tau_m,
 \qquad M_\Gamma K_{2m+1}=(2m+2)\tau_{m+1}.
\]
To prove this, start at $K_0=M_\Gamma^{-1}$.  Adding the odd term
in \eqref{lgamma:eq:terms} to the asserted even sum gives
$(2m+1)(2m+3/2)\tau_m/(2m+1/2)=(2m+2)\tau_{m+1}$.
Adding the next even term gives $(2m+3)\tau_{m+1}$.  This proves
\eqref{lgamma:eq:exactkernel} at every degree.

For the linear estimates, the exact original-coordinate coefficients are
\[
 b_n=\sqrt{\gamma_{n+1}/\gamma_n}=\sqrt{(n+1)(n+\tfrac12)},
 \qquad u_{n+1}=\frac{u_n+b_{n-1}u_{n-1}}{b_n}.
\]
For $n\ge1$,
\[
 b_n^2-(1+b_{n-1})^2
 =2n-\tfrac12-2\sqrt{n(n-\tfrac12)}>0,
\]
since both compared quantities are positive and
$(2n-1/2)^2-4n(n-1/2)=1/4$.  Therefore
\begin{equation}\label{lgamma:eq:pairmax}
 u_{n+1}\le\frac{1+b_{n-1}}{b_n}\max(u_{n-1},u_n)
 <\max(u_{n-1},u_n).
\end{equation}
The initial squares are $1/M_\Gamma$ and $2/M_\Gamma$, proving
\eqref{lgamma:eq:linear}.  The next squares are $3/(2M_\Gamma)$ and
$9/(5M_\Gamma)$.  Starting \eqref{lgamma:eq:pairmax} from this pair bounds
every term from index three on by $9/(5M_\Gamma)$.  Summing gives
$M_\Gamma K_N\le1+2+3/2+(N-2)9/5$, proving
\eqref{lgamma:eq:betterlinear}.
\end{proof}

\section{The disk estimate and the exact coefficient-to-observation map}

The imaginary-point domination is the argument in Christian Berg and
Ryszard Szwarc's Section 3 \cite{lgamma:BS}.  Their source starts with mass one
and $b_0=1$; the following proof retains the original finite mass and
recurrence coefficients.  Their label \texttt{thm:lambda} is a lemma,
not a theorem.

\begin{lgammalemma}\label{lgamma:lem:disk}
Let $m$ be a nonzero finite symmetric positive measure on $\lgammaR$, with
moments through order $2N$ and positive definite original monomial Gram
$H_y=(\int y^{j+\ell}\,dm)_{j,\ell=0}^N$.  If $P_n$ are its
orthonormal polynomials with positive leading coefficient and
$K_N^m(i,i)=\sum_{n=0}^N|P_n(i)|^2$, then
\begin{equation}\label{lgamma:eq:diskgram}
 |P_n(z)|\le|P_n(i)|\quad(|z|\le1),\qquad
 H_y\succeq K_N^m(i,i)^{-1}I.
\end{equation}
\end{lgammalemma}
\begin{proof}
Symmetry gives parity $P_n(-y)=(-1)^nP_n(y)$.  Expanding $yP_n$ in
the orthonormal basis shows that its coefficients below index $n-1$
vanish: pair instead $P_n$ with $yP_j$ of degree at most $n-1$.
The coefficient at index $n$ vanishes by parity.  The leading-coefficient
ratio gives $b_n>0$ at index $n+1$; symmetry of multiplication gives
$b_{n-1}$ at index $n-1$.  Consequently
\[
 yP_n=b_nP_{n+1}+b_{n-1}P_{n-1}.
\]
Starting at $P_0=m(\lgammaR)^{-1/2}$ and
$P_1(y)=m(\lgammaR)^{-1/2}y/b_0$, induction gives
$P_n(i)=i^nv_n$ with $v_n>0$ and
$v_{n+1}=(v_n+b_{n-1}v_{n-1})/b_n$.
The triangle inequality in the same recurrence proves disk domination
by induction.  Degree zero is immediate.

For $f(y)=\sum_{j=0}^Nc_jy^j=\sum_{n=0}^Nd_nP_n(y)$,
Cauchy--Schwarz and disk domination give on $|z|=1$
\[
 |f(z)|^2\le K_N^m(i,i)\sum|d_n|^2
 =K_N^m(i,i)c^*H_yc.
\]
Integrating over the circle gives exactly
$\sum|c_j|^2\le K_N^m(i,i)c^*H_yc$ by Fourier orthogonality.
No source mass or coordinate was changed.
\end{proof}

Let the original complex coordinate be $S=c+iy$, with the actual real
$c$ retained.  For $F(S)=\sum a_jS^j$ and $f(y)=F(c+iy)=\sum b_ry^r$,
binomial expansion gives $b=T_N(c)a$, with
\begin{equation}\label{lgamma:eq:T}
 T_{rj}=\binom jr c^{j-r}i^r,\qquad
 (T^{-1})_{rj}=\binom jr(-c)^{j-r}i^{-j}\quad(r\le j),
\end{equation}
and both entries zero if $r>j$.  The inverse formula follows by
expanding $((S-c)/i)^j$.  In particular,
\begin{equation}\label{lgamma:eq:Tgram}
 \det T=i^{N(N+1)/2},\quad H_S=T^*H_yT,\quad
 H_S^{-1}=T^{-1}H_y^{-1}T^{-*},\quad \det H_S=\det H_y.
\end{equation}
Here $H_S$ is the literal Gram of $1,S,\ldots,S^N$, and
$T^{-*}=(T^{-1})^*$.  Thus $H_y\succeq\beta I$ implies
$H_S\succeq\beta T^*T$, not an unsupported bound by $\beta I$.

Let $J_{S,N}$ be the actual native coefficient-to-class matrix and let
$\Lambda$ be the original observation onto its image in its fixed frame.
All physical-unit and divided-jet factors occurring in $J_{S,N}$ remain
there.  Put
\begin{equation}\label{lgamma:eq:A}
 A_{S,N}=\Lambda J_{S,N},\quad
 A_{y,N}=A_{S,N}T^{-1},\quad
 \mathcal B_N=A_{S,N}T^{-1}T^{-*}A_{S,N}^*.
\end{equation}
Whenever the source cutoff maps onto the class space, $A_{S,N}$ has full
row rank on the observation image, so $\mathcal B_N>0$.

\begin{lgammalemma}\label{lgamma:lem:min}
For a positive definite coefficient Gram $H$ and full-row-rank $A$,
\begin{equation}\label{lgamma:eq:min}
 \min_{Aa=w}a^*Ha=w^*(AH^{-1}A^*)^{-1}w,\qquad
 a_{\min}=H^{-1}A^*(AH^{-1}A^*)^{-1}w.
\end{equation}
For the unchanged Gamma source at $S=c+iy$ this gives
\begin{equation}\label{lgamma:eq:Qsigma}
 \boxed{Q_N^\sigma=(A_{S,N}(H_S^\sigma)^{-1}A_{S,N}^*)^{-1}
 \succeq\beta_N\mathcal B_N^{-1}.}
\end{equation}
\end{lgammalemma}
\begin{proof}
The proposed minimizer has image $w$.  Every other lift is $a_{\min}+h$
with $Ah=0$, and $Ha_{\min}=A^*(AH^{-1}A^*)^{-1}w$ is orthogonal
to $h$.  Its energy is the minimum in \eqref{lgamma:eq:min} plus $h^*Hh$.
Lemma \ref{lgamma:lem:disk} and Theorem \ref{lgamma:thm:kernel} give
$(H_y^\sigma)^{-1}\preceq\beta_N^{-1}I$.  Conjugate by
$A_{S,N}T^{-1}$ and invert to obtain \eqref{lgamma:eq:Qsigma}.  Inversion
reverses order because conjugating $0<X\preceq Y$ by $Y^{-1/2}$
gives a positive matrix with eigenvalues at most one, whose inverse has
eigenvalues at least one.
\end{proof}

For clarity, if $V_S,V_y$ are the literal derivative-jet matrices at
$s_\ell=c+iy_\ell$ with row orders $d$, the chain rule gives
\begin{equation}\label{lgamma:eq:jets}
 V_yT=D V_S,\qquad D_{(\ell,d),(\ell,d)}=i^d.
\end{equation}
It also holds for corresponding divided derivatives.  Thus if the
actual map is $J_{S,N}=U_{\rm phys}V_S$, its transform is
$J_{S,N}T^{-1}=U_{\rm phys}D^{-1}V_y$: the physical unit is not removed,
and $D^{-1}$ is not moved through it.  For square jet matrices,
\[
 \det V_S=i^{N(N+1)/2-\sum d}\det V_y.
\]
This keeps the determinant phase, including the case of zero determinants.

\section{The native simple-quartet domain, weight and exact window}

We now use exactly the simple-quartet domain of CAI and OB
\cite{lgamma:CAI,lgamma:OB}: a stipulated full simple quartet of zeros of $g=2\xi$,
\begin{equation}\label{lgamma:eq:nativedomain}
 \rho=\tfrac12+\delta+i\gamma,\quad 0<\delta<\tfrac12,\quad\gamma>2,
 \quad \mathcal R=\{\rho,\bar\rho,1-\rho,1-\bar\rho\},
 \quad h(s)=\prod_{\zeta\in\mathcal R}(s-\zeta).
\end{equation}
Here $\xi(s)=\tfrac12s(s-1)\pi^{-s/2}\Gamma(s/2)\zeta(s)$.
This domain does not assert existence of an off-line quartet.
All statements below concern the original objects on that stipulated
domain, without adding a zero, period or phase assumption.
Retain
\begin{equation}\label{lgamma:eq:nativeweight}
 v_h=g/h,\quad w_h(y)=\frac{|v_h(1/2+iy)|^2}{2\pi},\quad
 dm_k(y)=w_h^{*k}(y)\,dy,\quad
 \mu_h=\int w_h(y)\,dy.
\end{equation}
The arithmetic mass is $\mu_h^k$, not $M_\Gamma$ and not
$M_\Gamma^k$.  Fix the original period and its admitted observation map.
The original parameters and polynomial are
\begin{equation}\label{lgamma:eq:chik}
 \begin{gathered}
 k=4\ell+1\ge9,\quad c=k/2,\quad q=(k+1)^2,\quad S=c+iy,\\
 \chi_k(S)=\prod_{a,b=0}^k
 [S-c-(2a-k)\delta-i(2b-k)\gamma],\qquad E_k=\lgammaC[S]/(\chi_k).
 \end{gathered}
\end{equation}
The roots are distinct since $\delta,\gamma>0$.  The exact remainder map
is onto for $N\ge q-1$ and has kernel $\chi_k\lgammaC[S]_{\le N-q}$.
The complete original cutoff window is
\begin{equation}\label{lgamma:eq:window}
 q-1\le N\le2q,
 \qquad N_*=2q+2.
\end{equation}
The four return endpoints are $q-1,q,2q-1,2q$; the adjacent-current sum
uses $N=q-1,\ldots,2q-1$ and weights $1,2,\ldots,2,1$ of total $2q$.
The larger $N_*$ is exactly the CAI comparison degree used in OB13 and
FW7; it is not a replacement window.

\paragraph{The unit and the exact physical presentation.}
The original local Taylor unit is $\upsilon_h=j_hv_h$, invertible
because the full simple zeros have been divided out.  In the physical
tensor presentation the original map is
\[
 \eta:E_k\longrightarrow E_h^{\otimes k},\qquad
 \eta[P]=\upsilon_h^{\otimes k}[P(s_1+\cdots+s_k)],
 \qquad E_h=\lgammaC[s]/(h).
\]
It is injective: all grid centres in \eqref{lgamma:eq:chik} occur as sums of
the retained quartet.  For given $a,b$, choose $x=\max(0,a+b-k)$
and quartet multiplicities $x,a-x,b-x,k-a-b+x$; these are nonnegative
and have the required two marginal counts.  Vanishing at all tensor
points, after multiplying
by the invertible unit, is equivalent to divisibility by $\chi_k$.
Thus the actual physical coefficient-to-class map is
$J_{S,N}=\eta\circ\lgammaremm_{\chi_k}$, with codomain the original image
$\eta(E_k)$, whenever this is the native frame used.  Its observation
is the retained period map, whose included form is
$P_k\Pi^{\otimes k}$ on that image.  In CAI's literal remainder
presentation the identical composite is written
\[
 \jmath\Lambda_{\rm rem}
 =P_k\Pi^{\otimes k}U^{\otimes k}\iota,
 \qquad U=M_{j_hv_h},\qquad \iota[P]=[P(\textstyle\sum s_i)].
\]
These are equal maps, not a deletion of $U$.  If $R$ is the fixed
matrix of $\eta$ between the two retained class frames, their exact
relations are
\begin{equation}\label{lgamma:eq:unitframe}
 J_{S,N}=R J_{{\rm rem},N},\quad
 \Lambda=\Lambda_{\rm rem}R^{-1},\quad
 G_N=R^{-*}G_{{\rm rem},N}R^{-1},\quad
 a_i=R a_{i,\rm rem},\quad v=Rv_{\rm rem}.
\end{equation}
Consequently $A_{S,N}$ and every observed vector below are unchanged.
This retains the original unit, including its actual phase.  In
particular the unit matrix $\operatorname{diag}(a,\bar a,\bar a,a)$
of NC1--3 \cite{lgamma:NC} is not set to the identity.  Neither NC's
proper-source kernel nor its independent target minimum is substituted
for the present observation fibre.

\paragraph{The established CAI constants in the simple domain.}
The original constants from CAI16--18, with common multiplicity one,
are exactly
\begin{equation}\label{lgamma:eq:CAIconstants}
 \begin{gathered}
 \alpha=\pi/2,\quad p=7/2,\quad B_h=50,\quad M=25,\\
 \vartheta_h=\int_{-1}^1w_h(y)\,dy>0,\quad
 M_\alpha=\int_\lgammaR e^{\alpha y}w_h(y)\,dy\in(0,\infty),\\
 c_\Gamma=\sqrt2e^{-7/6},\qquad C_\Gamma=\sqrt{2\sqrt5}\,e^{2/3},\\
 a_k=\frac{c_h\vartheta_h^{k-3}e^{-\alpha(k-3)}(k-2)^{-50}}
              {C_\Gamma2^{25}},\qquad
 A_k=\frac{C_h^{\rm tilt}}{c_\Gamma}k^{9/2}M_\alpha^{k-1},\\
 D_n=[1+4(n+25)^2]^{25},\qquad
 \ell_k=\frac{a_k}{[1+4(2q+2+25)^2]^{25}},\qquad u_k=A_k.
 \end{gathered}
\end{equation}
Here $c_h$ is the original TW7 lower-envelope constant and
$C_h^{\rm tilt}$ the original BT12 compact/tail maximum as retained in
CAI16--17, not newly selected constants.  The established density
inequalities in that provider are, for every convolution order $b\ge3$,
\begin{align}
 w_h^{*b}(y)&\ge c_h\vartheta_h^{b-3}
 e^{-\alpha(|y|+b-3)}(b-2+|y|)^{-50},\label{lgamma:eq:CAIlowdens}\\
 w_h^{*b}(y)&\le bC_h^{\rm tilt}M_\alpha^{b-1}
 e^{-\alpha|y|}(1+|y|/b)^{-7/2},\label{lgamma:eq:CAIupdens}\\
 c_\Gamma e^{-\alpha|y|}(1+|y|)^{-1/2}
 &\le\frac{d\sigma}{dy}
 \le C_\Gamma e^{-\alpha|y|}(1+|y|)^{-1/2}.
 \label{lgamma:eq:CAIgammadens}
\end{align}
These quoted full-density results are the native analytic input.  For
completeness their passage to the exact polynomial receiver is as follows.
The elementary inequalities
$b-2+|y|\le(b-2)(1+|y|)$,
$(1+|y|)^{50}\le2^{25}(1+y^2)^{25}$, and
$1+|y|/b\ge(1+|y|)/b$ give
\[
 a_b(1+y^2)^{-25}\,d\sigma\le w_h^{*b}(y)\,dy\le A_b\,d\sigma.
\]
The original Gamma recurrence gives
$\lgammanorm{yP}_\sigma\le2(n+1)\lgammanorm P_\sigma$ at degree $n$:
its raising and lowering parts have respective norms at most $n+1$
and $n$ on that coefficient span.  Iterating at successive degrees gives
$\lgammanorm{y^jP}_\sigma\le[2(n+25)]^j\lgammanorm P_\sigma$ for $0\le j\le25$.
Expand $(1+y^2)^{25}$ to obtain
\[
 \int|P|^2(1+y^2)^{25}\,d\sigma\le D_n\lgammanorm P_\sigma^2.
\]
Cauchy--Schwarz with the reciprocal weights gives
$\int|P|^2(1+y^2)^{-25}\,d\sigma\ge D_n^{-1}\lgammanorm P_\sigma^2$.
Thus, for every complex polynomial of degree at most $N_*=2q+2$,
\begin{equation}\label{lgamma:eq:nativecomparison}
 \boxed{\ell_k\lgammanorm{P}_{\sigma,c}^{2}
       \le\lgammanorm{P}_{m_k,c}^{2}
       \le u_k\lgammanorm{P}_{\sigma,c}^{2}.}
\end{equation}
Here the subscript $c$ means that the same original polynomial is
evaluated at $S=c+iy$.  Its degree is unchanged by the explicitly
proved map $T_N(c)$.  This is CAI18 at $N_*$, exactly as used in
OB13 and OB18.  It does not identify $w_h^{*k}$ with $\sigma$ or with
$\sigma^{*k}$.

\section{Identical-fibre transfer to the native observed form}

Let $H_{S,N}^{\rm ar}$ be the full original coefficient Gram for
\eqref{lgamma:eq:nativeweight}, and let $G_N^{\rm ar}$ be the attained class
form at degree $N$, including every original relation.  In the retained
class frame,
\[
 (G_N^{\rm ar})^{-1}
 =J_{S,N}(H_{S,N}^{\rm ar})^{-1}J_{S,N}^*,\quad
 Q_N^{\rm ar}=(\Lambda(G_N^{\rm ar})^{-1}\Lambda^*)^{-1}.
\]
Composition therefore gives exactly
$Q_N^{\rm ar}=(A_{S,N}(H_{S,N}^{\rm ar})^{-1}A_{S,N}^*)^{-1}$.
All these matrices exist for \eqref{lgamma:eq:window}; the densities have
positive norm on nonzero polynomials, the remainder map is onto, and
the observation is taken onto its image.

\begin{lgammatheorem}[Native observed lower form]\label{lgamma:thm:nativelower}
At every original cutoff $q-1\le N\le2q$,
\begin{equation}\label{lgamma:eq:samefibre}
 \ell_kQ_N^\sigma\preceq Q_N^{\rm ar}\preceq u_kQ_N^\sigma,
\end{equation}
where the Gamma quotient uses the identical $A_{S,N}$, period,
physical-unit factors and observation frame.  Consequently
\begin{equation}\label{lgamma:eq:nativelower}
 \boxed{Q_N^{\rm ar}\succeq
 \ell_k\beta_N\left[
 \Lambda J_{S,N}T_N(k/2)^{-1}T_N(k/2)^{-*}J_{S,N}^*\Lambda^*
 \right]^{-1}.}
\end{equation}
The scalar $\beta_N$ is the exact formula \eqref{lgamma:eq:beta}.
\end{lgammatheorem}
\begin{proof}
Fix an arbitrary observation vector $w$.  The candidate set in both
minima is exactly
\[
 \{P\in\lgammaC[S]_{\le N}:A_{S,N}\operatorname{coeff}_S(P)=w\}.
\]
It is nonempty and is the same affine fibre, with every relation and
kernel lift included.  Apply \eqref{lgamma:eq:nativecomparison} to every
candidate.  Taking infima proves the lower inequality.  For the upper
inequality evaluate the arithmetic norm at the Gamma minimizer;
Lemma \ref{lgamma:lem:min} proves its existence and its exact norm.  This
proves \eqref{lgamma:eq:samefibre} for every $w$, hence in Loewner order.
Combine its lower inequality with \eqref{lgamma:eq:Qsigma} and retain
$\mathcal B_N$ from \eqref{lgamma:eq:A}.  This proves \eqref{lgamma:eq:nativelower}.
It is precisely the full-matrix version of OB18's same-class minimum
argument.  No extra large-period or conductor-corner hypothesis from
OB14 is needed or imposed.
\end{proof}

If the observation image has dimension zero, all observed vectors and
forms below are zero and the same assertions hold with empty-matrix
conventions.  Otherwise the inverses above are ordinary positive
definite inverses.

\section{The actual arithmetic boundary triple and its exact Gram}

At cutoff $N$ retain the \emph{arithmetic}, $S$-monic orthogonal
polynomials $p_n^{\rm ar}(S)$ and their full norms $\omega_n>0$.
As in FW1--5 and HG6--7 \cite{lgamma:FW,lgamma:HG}, use the original classes
\begin{equation}\label{lgamma:eq:nativetriple}
 a_1=b_N=[p_N^{\rm ar}],\qquad
 a_2=b_{N+1}=[p_{N+1}^{\rm ar}],\qquad
 v=v_\lambda=\left[\frac{\chi_k(S)}{(S-\lambda)\chi_k'(\lambda)}\right],
 \quad\lambda=k\rho.
\end{equation}
Brackets denote the original class, with the physical presentation
\eqref{lgamma:eq:unitframe} whenever that is the retained frame.  In particular
the complete physical unit multiplies these classes as specified there.
The vector $a_2$ is a class represented by a degree-$N+1$ polynomial,
but its minimum at cutoff $N$ includes all its degree-at-most-$N$
representatives; they exist because $N\ge q-1$.  It is not replaced
by a Gamma polynomial or by an inadmissible unreduced lift.

Write $G=G_N^{\rm ar}$, $Q=Q_N^{\rm ar}$, and
\begin{equation}\label{lgamma:eq:omegal}
 \Omega=\Lambda^*Q\Lambda,\qquad L=G-\Omega,\qquad
 X=(a_1,a_2,v),\qquad W=\Lambda X.
\end{equation}
The exact class-space minimum section is
$S_B=G^{-1}\Lambda^*Q$.  It satisfies $\Lambda S_B=I_B$.
Thus $P_K=I_E-S_B\Lambda$ maps onto $\ker\Lambda$, is the identity
there, and $I_K^*GS_B=0$ for the actual kernel inclusion
$I_K:\ker\Lambda\hookrightarrow E$.
The resulting orthogonal decomposition proves
\begin{equation}\label{lgamma:eq:split}
 0\preceq\Omega\preceq G,\quad L=GP_K\succeq0,
 \qquad X^*GX=X^*\Omega X+X^*LX.
\end{equation}
This is the exact FW37--38 and HR2 minimum, now proved in the actual
three columns.  Define all its entries without dropping any:
\begin{equation}\label{lgamma:eq:entries}
 \begin{gathered}
 d_i=a_i^*Ga_i,\quad E=v^*Gv,\quad F_i=a_i^*Gv,\\
 D_i=a_i^*\Omega a_i,\quad e=v^*\Omega v,\quad
 B_i=a_i^*\Omega v,\quad z=a_1^*\Omega a_2,\\
 k_i=a_i^*La_i=d_i-D_i,\quad
 K_i=a_i^*Lv=F_i-B_i,\quad e_K=v^*Lv=E-e.
 \end{gathered}
\end{equation}
The native source pairing $a_1^*Ga_2$ is zero, but the observed pairing
$z$ need not be zero.  To verify the former assertion in the original
objects, the exact functional equations $\xi(1-s)=\xi(s)$ and
$h(1-s)=h(s)$ make $w_h$ even, and convolution preserves evenness.
Evenness of $m_k$ makes $p_n^{\rm ar}(k-S)=(-1)^np_n^{\rm ar}(S)$.
Reflection permutes the full roots of $\chi_k$ and takes each entire
affine fibre to the reflected fibre isometrically.  Its induced map
$\Sigma$ therefore satisfies $\Sigma^*G\Sigma=G$, while
$\Sigma b_n=(-1)^nb_n$.  Hence
$a_1^*Ga_2=(-1)^{2N+1}a_1^*Ga_2=0$.  The same assertion in the
physical frame follows by the exact congruence \eqref{lgamma:eq:unitframe}.
Consequently the complete source, observed and kernel triple Grams are
\begin{equation}\label{lgamma:eq:threegrams}
 \begin{gathered}
 \mathcal F=X^*GX=
 \begin{pmatrix}d_1&0&F_1\\0&d_2&F_2\\\bar F_1&\bar F_2&E\end{pmatrix},\\
 \mathcal C=W^*QW=
 \begin{pmatrix}D_1&z&B_1\\\bar z&D_2&B_2\\\bar B_1&\bar B_2&e\end{pmatrix},\qquad
 \mathcal K=X^*LX=
 \begin{pmatrix}k_1&-z&K_1\\-\bar z&k_2&K_2\\\bar K_1&\bar K_2&e_K\end{pmatrix},\\
 \mathcal F=\mathcal C+\mathcal K,\qquad \mathcal C,\mathcal K\succeq0.
 \end{gathered}
\end{equation}
In particular the two complementary off-diagonal boundary terms are
$z$ and $-z$, not two separately vanishing terms.

\begin{lgammatheorem}[The joint native observed constraint]\label{lgamma:thm:joint}
Put
\begin{equation}\label{lgamma:eq:baseline}
 \mathfrak b_{k,N}=\ell_k\beta_N,\qquad
 \Psi_N=W^*\mathcal B_N^{-1}W,\qquad
 \mathcal C^{(0)}=\mathfrak b_{k,N}\Psi_N.
\end{equation}
Then the actual matrices in \eqref{lgamma:eq:threegrams} satisfy
\begin{equation}\label{lgamma:eq:joint}
 \boxed{\mathcal C\succeq\mathcal C^{(0)},\qquad
 \mathcal R:=\mathcal C-\mathcal C^{(0)}\succeq0,\qquad
 \mathcal F-\mathcal C^{(0)}=\mathcal K+\mathcal R,\quad
 \mathcal K\succeq0.}
\end{equation}
All entries of $\Psi_N$ use the same original $\Lambda a_1$,
$\Lambda a_2$, $\Lambda v$ and the complete original map in
\eqref{lgamma:eq:nativelower}.
Moreover $\ker\mathcal C=\ker\Psi_N=\ker W$; no independence of
the three observed columns is assumed.
\end{lgammatheorem}
\begin{proof}
Congruence of \eqref{lgamma:eq:nativelower} by the three-column matrix $W$
gives the first inequality and therefore the positive remainder.
The remaining equality follows from \eqref{lgamma:eq:threegrams}.  It
preserves the whole native kernel remainder rather than replacing it
with a trace or a rank.
Finally, for $t\in\lgammaC^3$, $t^*\mathcal Ct=(Wt)^*Q(Wt)$ vanishes
exactly when $Wt=0$, since $Q>0$.  The same argument with
$\mathcal B_N^{-1}>0$ proves the stated kernel identities.
\end{proof}

Here are all scalar positive-semidefinite constraints in explicit form.
Write $\psi_{ij}=(\Psi_N)_{ij}$ and set
\begin{equation}\label{lgamma:eq:resentries}
 \begin{gathered}
 x=D_1-\mathfrak b_{k,N}\psi_{11},\quad
 y=D_2-\mathfrak b_{k,N}\psi_{22},\quad
 h_B=e-\mathfrak b_{k,N}\psi_{33},\\
 \zeta_B=z-\mathfrak b_{k,N}\psi_{12},\qquad
 r_1=B_1-\mathfrak b_{k,N}\psi_{13},\quad
 r_2=B_2-\mathfrak b_{k,N}\psi_{23}.
 \end{gathered}
\end{equation}
Then $\mathcal R=\left(\begin{smallmatrix}
x&\zeta_B&r_1\\\bar\zeta_B&y&r_2\\\bar r_1&\bar r_2&h_B
\end{smallmatrix}\right)$, and \eqref{lgamma:eq:joint} is equivalent, as far
as its observed lower-form inequality is concerned, to the seven real
conditions
\begin{equation}\label{lgamma:eq:allminors}
 \boxed{\begin{gathered}
 x\ge0,\quad y\ge0,\quad h_B\ge0,\\
 xy-|\zeta_B|^2\ge0,\quad xh_B-|r_1|^2\ge0,
 \quad yh_B-|r_2|^2\ge0,\\
 xyh_B+2\Re(\zeta_Br_2\bar r_1)
 -x|r_2|^2-y|r_1|^2-h_B|\zeta_B|^2\ge0.
 \end{gathered}}
\end{equation}
The last line is the complete determinant, including its triple product.
Necessity follows from principal minors of a positive semidefinite
matrix.  For sufficiency, add $\epsilon I$ with $\epsilon>0$.
Its first and second leading principal minors are positive by the
first two lines.  Its determinant is
\[
 \det(\mathcal R+\epsilon I)
 =\det\mathcal R+
 \epsilon\sum_{|I|=2}\det\mathcal R_{I,I}
 +\epsilon^2\operatorname{tr}\mathcal R+\epsilon^3>0.
\]
Sylvester's criterion gives
$\mathcal R+\epsilon I>0$; let $\epsilon\downarrow0$.
Thus no positivity condition is missing at singular boundary cases.
These are exact constraints on the native matrix; they do not assert
that every numerical matrix satisfying them is realized by the native
arithmetic source.

For the un-subtracted observed Gram the corresponding explicit conditions
are
\begin{equation}\label{lgamma:eq:observedminors}
 \begin{gathered}
 D_1,D_2,e\ge0,\quad D_1D_2\ge|z|^2,\quad
 D_1e\ge|B_1|^2,\quad D_2e\ge|B_2|^2,\\
 D_1D_2e+2\Re(zB_2\bar B_1)
 -D_1|B_2|^2-D_2|B_1|^2-e|z|^2\ge0.
 \end{gathered}
\end{equation}
The kernel remainder is independently retained through
\begin{equation}\label{lgamma:eq:kernelminors}
 \begin{gathered}
 k_1,k_2,e_K\ge0,\quad k_1k_2\ge|z|^2,\quad
 k_1e_K\ge|K_1|^2,\quad k_2e_K\ge|K_2|^2,\\
 k_1k_2e_K-2\Re(zK_2\bar K_1)
 -k_1|K_2|^2-k_2|K_1|^2-e_K|z|^2\ge0.
 \end{gathered}
\end{equation}
The sign change in the last triple product is forced by the exact
off-diagonal $-z$ in \eqref{lgamma:eq:threegrams}.

For $h_B>0$, completing the square in the last coordinate of
$\mathcal R$ gives the equivalent two-dimensional Schur conditions
\begin{equation}\label{lgamma:eq:resdisk}
 x-\frac{|r_1|^2}{h_B}\ge0,\quad
 y-\frac{|r_2|^2}{h_B}\ge0,\qquad
 \left|\bar r_1r_2-h_B\bar\zeta_B\right|^2
 \le(h_Bx-|r_1|^2)(h_By-|r_2|^2).
\end{equation}
No phase was chosen: the disk has the actual complex centre
$h_B\bar\zeta_B$.  If $h_B=0$, \eqref{lgamma:eq:allminors} instead forces
$r_1=r_2=0$ and leaves the exact conditions $x,y\ge0$ and
$xy\ge|\zeta_B|^2$.  This covers the case without a division.
Similarly, when $e>0$ the un-subtracted constraint implies
\begin{equation}\label{lgamma:eq:Bdisk}
 |\bar B_1B_2-e\bar z|^2
 \le(eD_1-|B_1|^2)(eD_2-|B_2|^2).
\end{equation}
Here $e>0$ is equivalent to $\Lambda v\ne0$; when $e=0$, positivity
forces $B_1=B_2=0$.  No additional period locus is silently imposed.

\section{Both complementary cross products and every remainder term}

Retain the native scale $\omega_N>0$ from \eqref{lgamma:eq:nativetriple}.
The exact current formulas in HG7 \cite{lgamma:HG}, also obtained by
substituting HR5 into the original response products \cite{lgamma:HR}, are
\begin{equation}\label{lgamma:eq:currents}
 \begin{aligned}
 \Phi_{K,N}&=\frac2{\omega_N}\Re(\bar K_1K_2),\\
 \Phi_{B,N}&=\frac2{\omega_N}\Re(\bar B_1B_2),\\
 \Phi_{\times,N}&=\frac2{\omega_N}
        \Re(\bar K_1B_2+\bar B_1K_2).
 \end{aligned}
\end{equation}
For verification on the native nonzero-denominator domain, put
$\mathfrak c_N=\bar F_1F_2/(\omega_NE)$, $U=K_1/F_1$ and
$V=K_2/F_2$.  Multiplying $2E\mathfrak c_N$ respectively by
$\bar UV$, $(1-\bar U)(1-V)$, and
$\bar U(1-V)+(1-\bar U)V$ gives exactly
\eqref{lgamma:eq:currents}.  The undivided right sides give their exact
polynomial continuation across vanishing response denominators.

Put, without discarding the positive remainder in \eqref{lgamma:eq:joint},
\[
 B_i^0=\mathfrak b_{k,N}\psi_{i3},\qquad
 K_i^0=F_i-B_i^0,\qquad B_i=B_i^0+r_i,\quad K_i=K_i^0-r_i.
\]
Expansion of every product yields
\begin{equation}\label{lgamma:eq:currentremainders}
 \begin{aligned}
 \frac{\omega_N}2\Phi_{B,N}
 &=\Re\bigl(\bar B_1^0B_2^0+\bar B_1^0r_2
              +\bar r_1B_2^0+\bar r_1r_2\bigr),\\
 \frac{\omega_N}2\Phi_{K,N}
 &=\Re\bigl(\bar K_1^0K_2^0-\bar K_1^0r_2
              -\bar r_1K_2^0+\bar r_1r_2\bigr),\\
 \frac{\omega_N}2\Phi_{\times,N}
 &=\Re\bigl(\bar K_1^0B_2^0+\bar B_1^0K_2^0
       +(\bar K_1^0-\bar B_1^0)r_2\\
 &\hspace{30mm}+\bar r_1(K_2^0-B_2^0)-2\bar r_1r_2\bigr).
 \end{aligned}
\end{equation}
Every linear cross term and every quadratic remainder is displayed.
In their sum these remainder terms cancel algebraically, not by an
estimate, giving the unchanged total
\begin{equation}\label{lgamma:eq:currenttotal}
 \Phi_{K,N}+\Phi_{B,N}+\Phi_{\times,N}
 =\frac2{\omega_N}\Re(\bar F_1F_2)
 =2E\Re\mathfrak c_N.
\end{equation}
The source coefficient is retained in its exact full-source definition,
not replaced by a leading expression omitting its signed phase remainder.
The baseline $\mathcal C^{(0)}$ is a proved lower form, not a declaration
that $\mathcal R$ is small or zero.

For example, \eqref{lgamma:eq:resdisk} supplies the exact finite enclosure
\begin{equation}\label{lgamma:eq:currentdisk}
 \begin{split}
 \Big|\frac{\omega_N}2\Phi_{B,N}
 -\Re(\bar B_1^0B_2^0+\bar B_1^0r_2+\bar r_1B_2^0
                   +h_B\bar\zeta_B)\Big|
 \le\sqrt{(h_Bx-|r_1|^2)(h_By-|r_2|^2)}.
 \end{split}
\end{equation}
For $h_B=0$ the same formula holds because $r_1=r_2=0$.
This is a correlated constraint retaining the actual remainder entries,
not an assignment of unknown phases.  All identities and inequalities
hold at every cutoff in \eqref{lgamma:eq:window}; insertion into the original
current sum keeps its weights $1,2,\ldots,2,1$, and the original
four-endpoint return keeps $q-1,q,2q-1,2q$.

\section{The separate actual Gamma-convolution reference}

For completeness let $b\ge1$ be any integer.  The same recurrence
argument applies to the actual finite-measure convolution $\sigma^{*b}$,
but this is not the native weight in \eqref{lgamma:eq:nativeweight}.  Its exact
mass is $M_{\Gamma,b}=M_\Gamma^b$.  Fubini and \eqref{lgamma:eq:laplace} give
\[
 \int e^{sy}\,d\sigma^{*b}(y)=M_{\Gamma,b}(\cos s)^{-b/2}.
\]
For
\[
 \sum_{n\ge0}p_n^{[b]}(y)\frac{t^n}{n!}
 =(1+t^2)^{-b/4}e^{y\arctan t},
\]
the same two-generating-function calculation proves
$\gamma_n^{[b]}=M_{\Gamma,b}n!(b/2)_n$ and
\[
 p_{n+1}^{[b]}=yp_n^{[b]}-n(n-1+b/2)p_{n-1}^{[b]}.
\]
Write $\alpha_b=b/2$ and $b_n=\sqrt{(n+1)(n+\alpha_b)}$.
For $n\ge1$,
\[
 b_n^2-(1+b_{n-1})^2
 =2n+\alpha_b-1-2\sqrt{n(n+\alpha_b-1)}\ge0,
\]
because the first compared term is positive and its square minus the
second square is $(\alpha_b-1)^2$.  The positive imaginary-point
recurrence and the initial squares $M_{\Gamma,b}^{-1}$ and
$(\alpha_bM_{\Gamma,b})^{-1}$ therefore prove
\begin{equation}\label{lgamma:eq:conv}
 K_N^{\sigma^{*b}}(i,i)
 \le\frac{1+N\max(1,2/b)}{(\sqrt{2\pi})^b}.
\end{equation}
At $b=2$ the recurrence gives equal imaginary-point squares at every
index, and \eqref{lgamma:eq:conv} is an equality.  No assertion in the native
transfer used $m_k=\sigma^{*k}$.

\section{Exact finite verification and provenance}

The companion \texttt{verify\_exact.py} proves by exact rational
calculation the alternating ratios, closed kernel formula and linear
bounds through $N=512$; it checks original Gram norms, inverse traces
and Loewner differences through $N=10$.  It checks symbolic-$c$
coefficient inverses and determinant phases through $N=8$, a complete
complex original-$S$ jet quotient at $N=4$, and the separate
Gamma-convolution recurrence bounds for $1\le k\le12$, $N\le128$.
These finite checks supplement the all-degree proofs, not replace them.
All printed Gram factors retain their exact common scalar
$\sqrt{2\pi}$ symbolically.  The first original kernel values multiplied
by $\sqrt{2\pi}$ are $1,3,9/2,63/10,63/8,77/8,539/48$.
The separate \texttt{verify\_native\_gram\_identities.py} checks the full
Hermitian determinant, complex Schur-disk identity, singular-case
principal-minor criterion, native source--observed--kernel decomposition,
and all three current expansions as exact symbolic polynomial identities.
It assigns no numerical value to a native arithmetic coefficient or
period phase.

The source-use coverage for this derivation is precise: Berg--Szwarc's
original version-one \LaTeX, Section 3, the complete disk-domination
and kernel-eigenvalue lemma proofs and intervening proposition; CAI's
domain, original unit and observation definitions and CAI16--19;
OB's full supplied text, especially OB13 and OB18; HG and HR's full
supplied texts; FW1--7 and FW37--40; and NC's original unit/frame
definitions and NC24--27 distinguishing its separate proper-source
minimum.  No separate full-paper reading of Berg--Szwarc by this
independent derivation is claimed; the parent task read that source in
full.  No Gamma-polynomial theorem from CCM has been imported.

\begin{thebibliography}{9}
\bibitem{lgamma:BS}
Christian Berg and Ryszard Szwarc,
\emph{The smallest eigenvalue of Hankel matrices},
arXiv:0906.4506v1 (2009), Section 3, original-source labels
\texttt{thm:pni} and \texttt{thm:lambda}, both lemmas;
\url{https://arxiv.org/abs/0906.4506v1}.
The original author \LaTeX\ source is retained locally as
\path{sources/0906.4506v1/original_source/main.tex}.

\bibitem{lgamma:CAI}
Split-Zero research programme,
\emph{The complete original action relative to its arithmetic total},
complete native CAI provider reconstructed from the delivered continuation
of 16 September 2026, CAI1--3, CAI6--9 and CAI16--19.
Original provider: \path{providers/CAI_COMPLETE_PREREQUISITE.tex}.
The retained analytic constants are the source's TW7 and BT12 constants,
as recorded in its AMT4/CAI16--17.

\bibitem{lgamma:OB}
Split-Zero research programme,
\emph{The observed terminal class on the complete original window},
native provider \path{OBSERVED_CLASS_RETURN.tex}, OB1--2,
the degree-$2q+2$ comparison preceding OB13, and OB18--19.

\bibitem{lgamma:FW}
Split-Zero research programme,
\emph{Full-window source products and the retained class:
finite quotient comparisons in the original arithmetic metric},
19 September 2026, FW1--7 and FW37--40.
Original provider: \path{providers/FULL_WINDOW_SOURCE_PRODUCTS.tex}.

\bibitem{lgamma:HG}
Split-Zero research programme,
\emph{Subexponential heat cost and the original arithmetic currents},
19 September 2026, HG1 and HG6--11.
Original provider: \path{providers/HEAT_GROWTH_AND_ORIGINAL_CURRENTS.tex}.

\bibitem{lgamma:HR}
Split-Zero research programme,
\emph{Heat approximation errors for the original complex observation responses},
19 September 2026, HR2--5 and HR8--11.
Original provider: \path{providers/HEAT_OBSERVATION_RESPONSE_ERRORS.tex}.

\bibitem{lgamma:NC}
Split-Zero research programme,
\emph{Vanishing of the original proper-source kernel through conductor coprimality},
19 September 2026, NC1--3 and NC24--27.
Original provider: \path{providers/ORIGINAL_CONDUCTOR_COPRIMALITY.tex}.
Used here only to retain the actual unit/frame conventions and to avoid
substituting its separate proper-source minimum for the present fibre.
\end{thebibliography}

\endgroup

\clearpage\begingroup
\part*{Closed range of the original full theta source}
\newtheorem{lthetatheorem}{Theorem}
\newtheorem{lthetalemma}[lthetatheorem]{Lemma}
\newtheorem{lthetaproposition}[lthetatheorem]{Proposition}
\newtheorem{lthetacorollary}[lthetatheorem]{Corollary}
\newcommand{\lthetaB}{\mathscr B}
\newcommand{\lthetaSch}{\mathcal S}
\newcommand{\lthetaHH}{\mathcal H}
\newcommand{\lthetaZ}{\mathbb Z}
\newcommand{\lthetaR}{\mathbb R}
\newcommand{\lthetaC}{\mathbb C}
\newcommand{\lthetaN}{\mathbb N}
\newcommand{\lthetaid}{\operatorname{id}}
\newcommand{\lthetacl}{\operatorname{cl}}
\newcommand{\lthetasupp}{\operatorname{supp}}
\newcommand{\lthetaF}{\mathcal F}



\section{The fixed objects and the source theorem}

Every space and map in this note is over $\lthetaC$. The original source is
\[
 V=\left\{\phi\in\lthetaSch(\lthetaR):\phi(-x)=\phi(x),\quad
       \phi(0)=0,\quad \int_\lthetaR\phi(x)\,dx=0\right\},
\]
with the subspace topology from the ordinary Schwartz space. Put
\[
 D=-x\frac{d}{dx},\qquad
 p_{b,j}(F)=\sup_{x>0}x^b|D^jF(x)|\quad(b\in\lthetaZ,\ j\in\lthetaN_0),
\]
\[
 \lthetaB=\{F\in C^\infty(\lthetaR_{>0}):p_{b,j}(F)<\infty
                     \text{ for every }b,j\}.
\]
The topology on $\lthetaB$ is exactly that generated by all $p_{b,j}$. The
operators and Fourier convention are
\begin{equation}\label{ltheta:eq:original}
 \Theta\phi(x)=2\sum_{n=1}^{\infty}\phi(nx),\qquad
 JF(x)=x^{-1}F(x^{-1}),\qquad
 \widehat\phi(\xi)=\int_\lthetaR\phi(x)e^{-2\pi i x\xi}\,dx.
\end{equation}
No conjugation by $x^{1/2}$, renormalization of $\Theta$, change of the
coordinate $x$, or substitution of a Hilbert norm is made here.

The literature source is Ralf Meyer, \emph{A spectral interpretation for
the zeros of the Riemann zeta function}, arXiv:math/0412277v3,
\href{https://arxiv.org/abs/math/0412277v3}{original versioned source}.
Its Section 2 defines
\[
 \lthetaHH_- =\bigcap_{s\in\lthetaR}\lthetaSch(\lthetaR_+^*)_s,\qquad
 \lthetaSch(\lthetaR_+^*)_s=
 \{F:x\mapsto x^sF(x)\text{ is Schwartz in }\log x\}.
\]
Its Section 3 defines $Z\phi(x)=\sum_{n\geq1}\phi(nx)$ and
$\lthetaHH_\cap=\{\phi\in\lthetaSch(\lthetaR):\phi\text{ even},\,
\phi(0)=\lthetaF\phi(0)=0\}$. Theorem 3.3, with source label
\texttt{the:Zeta\_estimate}, proves closed range and the
topological-isomorphism-onto-range property, including the restriction
$Z:\lthetaHH_\cap\to\lthetaHH_-$. Its proof uses Poisson summation and
Proposition 3.2 (label \texttt{pro:Z\_invertible}), together with
Proposition 2.1 (label \texttt{pro:Sch\_estimates}).

Meyer's Fourier kernel has the opposite sign. For an even $\phi$ the
change of variable $x\mapsto-x$ in the defining integral proves that
the two transforms are equal, as functions, with no scalar factor.
Consequently $V=\lthetaHH_\cap$ with identical source topology. His $Z$ and
the original map satisfy $\Theta=2Z$. The next sections prove the target
topology identification and independently reconstruct the closed-range
conclusion with explicit inverse estimates.

\section{Identity equivalence of the complete seminorm systems}

For $s\in\lthetaR$ and $k,j\in\lthetaN_0$ define
\begin{equation}\label{ltheta:eq:q}
 q_{s,k,j}(F)=\sup_{x>0}(1+|\log x|)^k
 \left|\left(x\frac d{dx}\right)^j(x^sF(x))\right|.
\end{equation}
These are the ordinary Schwartz seminorms in the variable $\log x$
after multiplication by $x^s$, expressed on the original $x$-axis. Their
projective family generates precisely Meyer's topology on $\lthetaHH_-$.

\begin{lthetatheorem}\label{ltheta:thm:topology}
The identity on positive-axis functions is a topological linear
isomorphism $\lthetaB=\lthetaHH_-$. More explicitly, set
\[
 n_-(s)=\lfloor s\rfloor-1,\qquad n_+(s)=\lceil s\rceil+1,
 \qquad C_k=\sup_{t\geq0}(1+t)^ke^{-t}.
\]
Here $C_0=1$ and $C_k=k^ke^{1-k}$ for integer $k\geq1$.
For every real $s$ and every $k,j\geq0$,
\begin{equation}\label{ltheta:eq:q-from-p}
 q_{s,k,j}(F)\leq C_k\sum_{\ell=0}^j\binom j\ell
 |s|^{j-\ell}\bigl(p_{n_-(s),\ell}(F)+p_{n_+(s),\ell}(F)\bigr).
\end{equation}
Conversely, for every integer $b$ and $j\geq0$,
\begin{equation}\label{ltheta:eq:p-from-q}
 p_{b,j}(F)\leq\sum_{\ell=0}^j\binom j\ell
 |b|^{j-\ell}q_{b,0,\ell}(F).
\end{equation}
In these formulas a zeroth power, including $0^0$, means $1$.
\end{lthetatheorem}

\begin{proof}
Write $L=x\,d/dx=-D$. The Leibniz rule, iterated $j$ times, gives the
exact operator identity
\[
 L^j(x^sF)=x^s(s-D)^jF
   =x^s\sum_{\ell=0}^j\binom j\ell s^{j-\ell}(-1)^\ell D^\ell F.
\]
For $x\geq1$ put $t=\log x$. Since $n_+(s)-s\geq1$,
\[
 (1+|\log x|)^kx^s\leq C_kx^{n_+(s)}.
\]
For $0<x\leq1$ put $t=-\log x$. Since $s-n_-(s)\geq1$,
\[
 (1+|\log x|)^kx^s\leq C_kx^{n_-(s)}.
\]
Substitution in the displayed binomial expansion proves
\eqref{ltheta:eq:q-from-p}. The constant follows by maximizing
$k\log(1+t)-t$ on $t\geq0$.

The conjugate identity $(b-L)(x^bF)=x^bDF$ yields
\[
 x^bD^jF=(b-L)^j(x^bF).
\]
Its binomial expansion gives \eqref{ltheta:eq:p-from-q}. Hence finiteness of
the full $p$-family and of the full $q$-family is equivalent, and each
seminorm of either family is bounded by a finite linear combination
from the other. This proves equality of the sets and of their locally
convex topologies, through the identity map.
\end{proof}

For completeness, $\lthetaB$ is a Fr\'echet space directly in the $p$-topology.
The seminorm family is countable and separates points. If $(F_n)$ is
Cauchy in every $p_{b,j}$, it is Cauchy in $C^\infty(K)$ for each compact
$K\subset(0,\infty)$, because
\begin{equation}\label{ltheta:eq:R}
 x^m\partial_x^m=R_m(D),\qquad
 R_m(t)=(-1)^m\prod_{r=0}^{m-1}(t+r),\qquad R_0(t)=1.
\end{equation}
The identity follows by induction, using
$x^{m+1}\partial^{m+1}=(x\partial-m)x^m\partial^m$.
Local limits therefore give a single smooth $F$. For a fixed $b,j$,
the $p_{b,j}$-Cauchy bound on $F_n-F_m$ passes pointwise to $m\to\infty$
and then to the supremum; thus $p_{b,j}(F_n-F)\to0$ and $F\in\lthetaB$.
The same local-limit argument proves completeness of $\lthetaSch(\lthetaR)$.
Parity, evaluation at zero, and integration are continuous Schwartz
functionals/conditions (for integration use the seminorm with weight
$(1+|x|)^2$). Therefore $V$ is a closed, complete Fr\'echet subspace.

\section{Continuity, Poisson symmetry, and exact M\"obius inversion}

Write
\[
 T_{N,j}(\psi)=\sup_{t\geq1}t^N|D^j\psi(t)|
 \qquad(\psi\in\lthetaSch(\lthetaR)).
\]
Each $T_{N,j}$ is a continuous Schwartz seminorm: expanding
$D^j=(-x\partial_x)^j$ gives a finite sum of polynomial coefficients
times ordinary derivatives. For $N>1$, termwise differentiation of
the theta series on every compact subset of $(0,\infty)$ is justified
by the convergent majorant $C n^{-N}$. On $x\geq1$ it gives
\begin{equation}\label{ltheta:eq:theta-tail}
 |D^j\Theta\psi(x)|\leq2\zeta(N)x^{-N}T_{N,j}(\psi).
\end{equation}

The ordinary Fourier transform maps $\lthetaSch(\lthetaR)$ continuously to itself.
Indeed, its differentiated, polynomially weighted expressions are
Fourier transforms of finite sums of derivatives of polynomially
weighted $\psi$, multiplied by the explicit powers of $2\pi i$ from
the kernel in \eqref{ltheta:eq:original}; their suprema are bounded by the
corresponding $L^1$ norms. Those norms are bounded by finitely many
Schwartz seminorms, using the integrable weight $(1+|x|)^{-2}$.
Fourier inversion gives $\widehat{\widehat\psi}(x)=\psi(-x)$.
It follows that the Fourier transform is an involution of $V$:
evenness is preserved, $\widehat\phi(0)=\int\phi=0$, and
$\int\widehat\phi=\phi(0)=0$.

Poisson summation with the exact kernel \eqref{ltheta:eq:original} says
\[
 \sum_{n\in\lthetaZ}\psi(nx)
    =x^{-1}\sum_{n\in\lthetaZ}\widehat\psi(n/x).
\]
It follows, for example, by periodizing $t\mapsto\psi(xt)$:
its $k$th Fourier-series coefficient is $x^{-1}\widehat\psi(k/x)$,
and the series and all needed derivatives converge absolutely by
Schwartz decay; evaluation at $t=0$ gives the formula.
For even $\psi$ this is exactly
\begin{equation}\label{ltheta:eq:poisson}
 \psi(0)+\Theta\psi(x)
       =x^{-1}\widehat\psi(0)+J\Theta\widehat\psi(x).
\end{equation}
Both endpoint terms vanish for $\phi\in V$, so
\begin{equation}\label{ltheta:eq:symmetry}
 \Theta\phi=J\Theta\widehat\phi,
 \qquad J\Theta\phi=\Theta\widehat\phi.
\end{equation}
The second identity uses $J^2=\lthetaid$.

Direct differentiation on the original axis gives
\begin{equation}\label{ltheta:eq:J}
 D^jJ=J(1-D)^j,\qquad
 p_{b,j}(JF)\leq\sum_{r=0}^j\binom jr p_{1-b,r}(F).
\end{equation}
For $j=1$ the first identity follows from
$-x\partial_x(x^{-1}F(x^{-1}))
=x^{-1}F(x^{-1})+x^{-2}F'(x^{-1})$;
induction proves the general case. In the seminorm bound substitute
$y=x^{-1}$, so $x^b x^{-1}=y^{1-b}$, and expand $(1-D)^j$.
In particular $J$ is a continuous involution of $\lthetaB$.

Choose the integers $N=\max(2,b)$ and $N'=\max(2,1-b)$.
Combining \eqref{ltheta:eq:theta-tail} on $x\geq1$ with
\eqref{ltheta:eq:symmetry} and \eqref{ltheta:eq:J} on $0<x\leq1$ yields
\begin{equation}\label{ltheta:eq:theta-continuity}
 p_{b,j}(\Theta\phi)
 \leq2\zeta(N)T_{N,j}(\phi)
  +2\zeta(N')\sum_{r=0}^j\binom jrT_{N',r}(\widehat\phi).
\end{equation}
Thus $\Theta:V\to\lthetaB$ is well defined and continuous in the precise
original topologies.

For a positive-axis smooth function define
$p^+_{u,j}(\psi)=\sup_{x>0}x^u|D^j\psi(x)|$ whenever finite;
the superscript distinguishes a source estimate from the target
seminorm notation. Let $\mu(n)$ be the M\"obius function and put
\[
 (MG)(x)=\sum_{n=1}^{\infty}\mu(n)G(nx).
\]
Whenever $p^+_{u,r}(G)<\infty$ for $u>1$ and $0\leq r\leq j$, its
series may be differentiated $j$ times locally uniformly and
\begin{equation}\label{ltheta:eq:mobius-bound}
 p^+_{u,j}(MG)\leq\zeta(u)p^+_{u,j}(G).
\end{equation}
Here $|\mu(n)|\leq1$ and $D^j[G(nx)]=(D^jG)(nx)$ give the estimate
term by term. In particular this applies to every $G\in\lthetaB$ and
integer $u\geq2$.

For $\phi\in V$ and fixed $x>0$, the double series
\[
 \sum_{n,m\geq1}|\mu(n)\phi(mnx)|
 \leq x^{-u}p^+_{u,0}(\phi)\zeta(u)^2
 \qquad(u>1)
\]
is finite. Grouping by $k=mn$ and using
$\sum_{n\mid k}\mu(n)=1$ for $k=1$ and $0$ otherwise gives
\begin{equation}\label{ltheta:eq:mobius-inversion}
 M\Theta\phi=2\phi\quad\text{on }(0,\infty).
\end{equation}
The divisor identity follows by taking the product
$\prod_{p\mid k}(1-1)$ when $k>1$; for $k=1$ the empty product is $1$.
Thus $\Theta$ is injective. Writing $F=\Theta\phi$, the exact inverse
formulas and their estimates are
\begin{align}
 \phi(x)&=\tfrac12\sum_{n\geq1}\mu(n)F(nx),&
 p^+_{u,j}(\phi)&\leq\tfrac12\zeta(u)p_{u,j}(F),
       \label{ltheta:eq:inverse-phi}\\
 \widehat\phi(x)&=\tfrac12\sum_{n\geq1}\mu(n)(JF)(nx),&
 p^+_{u,j}(\widehat\phi)&\leq\tfrac12\zeta(u)
                    \sum_{\ell=0}^j\binom j\ell p_{1-u,\ell}(F),
       \label{ltheta:eq:inverse-hat}
\end{align}
for $x>0$, integers $u\geq2$, and $j\geq0$. The second line uses
\eqref{ltheta:eq:symmetry}. These are inverse formulas on the image; they
do not assert that $MG$ extends to an element of $V$ for an arbitrary
$G\in\lthetaB$.

\section{Explicit recovery of every source Schwartz seminorm}

The following estimate supplies the part of the inverse-continuity
argument where both the function and its Fourier transform are needed.
All $L^2$ norms in this section use ordinary Lebesgue measure $dx$ or
$d\xi$ on $\lthetaR$, and the Fourier transform is exactly
\eqref{ltheta:eq:original}. We use its Plancherel norm equality.

\begin{lthetalemma}\label{ltheta:lem:uncertainty}
For $h\in\lthetaSch(\lthetaR)$ and integers $M,N\geq0$,
\begin{equation}\label{ltheta:eq:uncertainty}
 \|h\|_2\leq2\left(4^M\|x^Mh\|_2
                         +4^N\|\xi^N\widehat h\|_2\right).
\end{equation}
\end{lthetalemma}

\begin{proof}
Let $P$ and $Q$ be multiplication by the indicator of $[-1/4,1/4]$
on the physical and Fourier sides, respectively. The operator
$P\lthetaF^{-1}Q$ has kernel $1_{[-1/4,1/4]}(x)
e^{2\pi i x\xi}1_{[-1/4,1/4]}(\xi)$; its Hilbert--Schmidt norm is
$\sqrt{(1/2)(1/2)}=1/2$. Therefore, by Plancherel and the triangle
inequality,
\[
 \|h\|_2\leq\|Ph\|_2+\|(1-P)h\|_2
 \leq\tfrac12\|h\|_2+\|(1-Q)\widehat h\|_2+\|(1-P)h\|_2.
\]
On $|x|>1/4$, $1\leq4^M|x|^M$; the analogous statement holds in
frequency. Moving the half-norm to the left proves the assertion.
\end{proof}

Write the exact polynomial from \eqref{ltheta:eq:R} as
$R_m(t)=\sum_{j=0}^m r_{m,j}t^j$. No coefficient is dropped in what
follows. For integers $c\geq3$ and $m\geq0$ put
\begin{align}
 E_{c,m}(F)&=\frac1{\sqrt2}\sum_{j=0}^m|r_{m,j}|
 \left(\frac{\zeta(c-1)}{\sqrt3}p_{c-1,j}(F)
                      +\zeta(c+1)p_{c+1,j}(F)\right),
                     \label{ltheta:eq:E}\\
 E^J_{c,m}(F)&=\frac1{\sqrt2}\sum_{j=0}^m|r_{m,j}|
 \left(\frac{\zeta(c-1)}{\sqrt3}
                     \sum_{\ell=0}^j\binom j\ell p_{2-c,\ell}(F)
       +\zeta(c+1)\sum_{\ell=0}^j\binom j\ell p_{-c,\ell}(F)\right).
                     \label{ltheta:eq:EJ}
\end{align}
For $a,b\in\lthetaN_0$ set, with $(b)_r=b!/(b-r)!$,
\begin{multline}\label{ltheta:eq:A}
 A_{a,b}(F)=2\left[4^{b+3}E_{a+3,b}(F)\right.\\
 \left.{}+4^{a+3}(2\pi)^{b-a}
 \sum_{r=0}^{\min(a,b)}\binom ar(b)_rE^J_{b+3,a-r}(F)\right].
\end{multline}
Every $A_{a,b}$ is explicitly a finite nonnegative linear combination
of original $p_{b,j}$ seminorms.

\begin{lthetaproposition}\label{ltheta:prop:inverse}
For $\phi\in V$, $F=\Theta\phi$, and $a,b\in\lthetaN_0$,
\begin{equation}\label{ltheta:eq:L2-inverse}
 \|x^a\phi^{(b)}\|_2\leq A_{a,b}(F).
\end{equation}
In particular the ordinary Schwartz seminorms satisfy
\begin{equation}\label{ltheta:eq:sup-inverse}
 \sup_{x\in\lthetaR}|x^a\phi^{(b)}(x)|
 \leq A_{a,b}(F)+aA_{a-1,b}(F)+A_{a,b+1}(F),
\end{equation}
where the middle term is omitted when $a=0$.
\end{lthetaproposition}

\begin{proof}
For any even Schwartz $\psi$, splitting the positive-axis integral at
$1$ gives
\[
 \|x^cD^j\psi\|_{L^2(\lthetaR)}
 \leq\sqrt2\left(\frac{p^+_{c-1,j}(\psi)}{\sqrt3}
                               +p^+_{c+1,j}(\psi)\right).
\]
Indeed, on $(0,1)$ the integrand's absolute value is at most
$xp^+_{c-1,j}$, whose squared integral is $(p^+_{c-1,j})^2/3$;
on $(1,\infty)$ it is at most $x^{-1}p^+_{c+1,j}$, whose squared
integral is $(p^+_{c+1,j})^2$. Evenness supplies the factor $\sqrt2$;
the square root of the sum is at most the sum of square roots.
Using the full polynomial $R_m$ and then
\eqref{ltheta:eq:inverse-phi}--\eqref{ltheta:eq:inverse-hat}, with $c-1\geq2$, proves
\begin{equation}\label{ltheta:eq:U-bound}
 \|x^cR_m(D)\phi\|_2\leq E_{c,m}(F),\qquad
 \|x^cR_m(D)\widehat\phi\|_2\leq E^J_{c,m}(F).
\end{equation}

Apply Lemma~\ref{ltheta:lem:uncertainty} to
$h=x^a\phi^{(b)}$, with $M=b+3$ and $N=a+3$. The first weighted
term is, by \eqref{ltheta:eq:R},
\[
 x^Mh=x^{a+b+3}\phi^{(b)}=x^{a+3}R_b(D)\phi.
\]
The exact Fourier identities for the specified negative-sign kernel are
\[
 \widehat{\partial_x^b\phi}(\xi)=(2\pi i\xi)^b\widehat\phi(\xi),
 \qquad
 \widehat{x^a\psi}(\xi)=\left(-\frac1{2\pi i}\right)^a
                                      \partial_\xi^a\widehat\psi(\xi).
\]
Their Leibniz expansion therefore gives
\[
 \xi^{a+3}\widehat h(\xi)
 =\left(-\frac1{2\pi i}\right)^a(2\pi i)^b
 \sum_{r=0}^{\min(a,b)}\binom ar(b)_r
       \xi^{a+3+b-r}\widehat\phi^{(a-r)}(\xi).
\]
For each summand the final factor is exactly
$\xi^{b+3}R_{a-r}(D)\widehat\phi(\xi)$, where
$D=-\xi\partial_\xi$ on this axis. The scalar prefactor has modulus
$(2\pi)^{b-a}$. Substituting \eqref{ltheta:eq:U-bound} into
\eqref{ltheta:eq:uncertainty} proves \eqref{ltheta:eq:L2-inverse} with exactly
\eqref{ltheta:eq:A}.

For a Schwartz function $h$, integration from $-\infty$ to $x$ gives
$|h(x)|^2\leq2\|h\|_2\|h'\|_2$, whence
$\|h\|_\infty\leq\|h\|_2+\|h'\|_2$.
For $h=x^a\phi^{(b)}$ its derivative is
$a x^{a-1}\phi^{(b)}+x^a\phi^{(b+1)}$. The preceding $L^2$ estimate
and the triangle inequality yield \eqref{ltheta:eq:sup-inverse}. The
seminorms on its left generate the full Schwartz topology; weights
$(1+|x|)^a$ are bounded by the finite sum
$\sum_{r=0}^a\binom ar|x|^r$ and conversely each $|x|^a$ is bounded by
$(1+|x|)^a$.
\end{proof}

\section{Closed image, the exact quotient, and the chain isomorphism}

\begin{lthetatheorem}\label{ltheta:thm:closed}
The original map $\Theta:V\to\lthetaB$ is a continuous linear injection,
a topological isomorphism onto its range with the subspace topology,
and its range is closed. In particular
\begin{equation}\label{ltheta:eq:closed}
 \overline{\Theta V}^{\,\lthetaB}=\Theta V.
\end{equation}
\end{lthetatheorem}

\begin{proof}
Continuity is \eqref{ltheta:eq:theta-continuity}, injectivity is
\eqref{ltheta:eq:mobius-inversion}, and inverse continuity is
\eqref{ltheta:eq:sup-inverse}. Suppose $F_n=\Theta\phi_n$ converges in $\lthetaB$
to $F$. Applying \eqref{ltheta:eq:sup-inverse} to
$\phi_n-\phi_m\in V$ shows that $(\phi_n)$ is Cauchy for every
Schwartz seminorm. Completeness and closedness of $V$ give
$\phi_n\to\phi\in V$. Continuity of $\Theta$ implies
$\Theta\phi_n\to\Theta\phi$ in $\lthetaB$, and its Hausdorff property gives
$F=\Theta\phi$. Thus the range is sequentially closed. Since $\lthetaB$ is
metrizable by its countable seminorm family, it is closed.
\end{proof}

The literature application reaches the same conclusion immediately
after Theorem~\ref{ltheta:thm:topology}: Meyer's restricted closed embedding
is precisely $Z=\Theta/2$ with source $V=\lthetaHH_\cap$ and target
$\lthetaHH_-=\lthetaB$. Since $2V=V$, its image equals $\Theta V$ as a subset of
that identical target. The independent estimates above additionally
make the inverse continuity completely explicit in original seminorms.

Let $Q=\lthetaB/\Theta V$ with its quotient topology. It is Hausdorff:
if a coset is nonzero, its representative lies outside the closed
subspace $\Theta V$, and the quotient separates it from the zero
coset. It is also Fr\'echet. One direct completeness argument chooses
an increasing defining seminorm sequence for $\lthetaB$, uses the quotient
seminorms, selects from a quotient-Cauchy sequence a subsequence with
successive quotient distances below $2^{-n}$ in each of the first
$n$ seminorms, and lifts each successive difference to a vector whose
same finitely many seminorms are below $2^{1-n}$. The lifted series is
Cauchy in each seminorm, hence converges in $\lthetaB$ and projects to the
subsequence limit; the original Cauchy sequence has the same limit.
The quotient is metrizable by these countably many quotient seminorms.

The original algebraic exact sequence labelled (E5),
\[
 0\longrightarrow
 \overline{\Theta V}^{\,\lthetaB}/\Theta V\longrightarrow Q
 \longrightarrow\lthetaB/\overline{\Theta V}^{\,\lthetaB}\longrightarrow0,
\]
therefore has its first nontrivial term equal to zero. Its last map
is exactly
\begin{equation}\label{ltheta:eq:E5}
 Q\xrightarrow{\ \cong\ }\lthetaB/\overline{\Theta V}^{\,\lthetaB},
 \qquad [F]_{\Theta V}\longmapsto[F]_{\overline{\Theta V}^{\,\lthetaB}},
\end{equation}
with inverse the same representative map. Both maps are continuous
because the two equivalence relations and quotient topologies are
identical. No new completion or representative change is involved.

To retain the operator factor at the chain level, define two-term
complexes in degrees $0,1$ by
\[
 C_\Theta=[V\xrightarrow{\Theta}\lthetaB],\qquad
 C_Z=[\lthetaHH_\cap\xrightarrow{Z}\lthetaHH_-].
\]
The explicitly typed chain isomorphism is
\begin{equation}\label{ltheta:eq:chain}
 T^0=2\lthetaid:V\longrightarrow\lthetaHH_\cap,\qquad
 T^1=\lthetaid:\lthetaB\longrightarrow\lthetaHH_-.
\end{equation}
Indeed $ZT^0\phi=Z(2\phi)=\Theta\phi=T^1\Theta\phi$.
Its inverse has degree-zero map $\tfrac12\lthetaid$ and degree-one map
$\lthetaid$. Thus the induced degree-one quotient isomorphism is
$[F]\mapsto[F]$, and degree-zero cohomology is zero on both sides
by injectivity. These are the range complexes associated to the
operators; no different complex is attributed to Meyer.

For the original two-leg differential the equally literal diagram is
\[
 [V\oplus V\xrightarrow{\,\Theta-J\Theta\,}\lthetaB]
 \xrightarrow{\ (2\lthetaid\oplus2\lthetaid,\,\lthetaid)\ }
 [\lthetaHH_\cap\oplus\lthetaHH_\cap\xrightarrow{\,Z-JZ\,}\lthetaHH_-].
\]
The equality of differentials follows separately for each signed leg.
Its inverse scales both source coordinates by $1/2$ and fixes the
target. Since $J\Theta=\Theta\lthetaF$ on $V$, the source map
$(\phi_+,\phi_-)\mapsto\phi_+-\widehat\phi_-$ and the target identity
define a surjective chain map from the original two-leg complex to
$C_\Theta$. Its kernel is the degree-zero graph
$\{(\widehat\psi,\psi):\psi\in V\}$ with zero degree-one term:
both inclusions follow directly from
$\phi_+-\widehat\phi_-=0$. In particular the two-leg degree-one image
is also precisely $\Theta V$, not a different closure.

Finally $D(V)\subset V$: it preserves evenness, $D\phi(0)=0$, and
integration by parts gives $\int_\lthetaR D\phi=\int_\lthetaR\phi=0$.
It acts continuously on $\lthetaB$ because $p_{b,j}(DF)=p_{b,j+1}(F)$,
and termwise differentiation gives $D\Theta=\Theta D$. Thus the
one-leg chain isomorphism commutes with $D$ in both degrees.
The signs for the two-leg complex are also explicit. Integration by
parts and differentiation of the Fourier kernel give
\[
 \lthetaF D\phi=(1+\xi\partial_\xi)\widehat\phi
           =(1-D)\lthetaF\phi,
 \qquad \lthetaF(1-D)\phi=D\lthetaF\phi.
\]
Consequently its source generator is $D\oplus(1-D)$ and its target
generator is $D$: by \eqref{ltheta:eq:J},
\[
 D(\Theta\phi_+-J\Theta\phi_-)
    =\Theta D\phi_+-J\Theta(1-D)\phi_-.
\]
The two-leg scaling isomorphism commutes with these generators.
So does the projection to the one-leg complex, because
$\phi_+-\widehat\phi_-$ is sent under the source generator to
$D\phi_+-\lthetaF(1-D)\phi_-=D(\phi_+-\widehat\phi_-)$.
No common diagonal generator on the two source legs and no sign
change to $x\partial_x$ is silently substituted.

\section{The exact scope for subspaces, finite cutoffs, and Hilbert norms}

\begin{lthetaproposition}\label{ltheta:prop:subspace}
For every linear subspace $W\subseteq V$, with closure taken in the
indicated original spaces,
\begin{equation}\label{ltheta:eq:W}
 \overline{\Theta W}^{\,\lthetaB}
       =\Theta\left(\overline W^{\,V}\right).
\end{equation}
The restriction of $\Theta$ induces a topological linear isomorphism
\[
 \overline W^{\,V}/W
       \longrightarrow\overline{\Theta W}^{\,\lthetaB}/\Theta W,
 \qquad[\psi]\longmapsto[\Theta\psi],
\]
with the respective quotient topologies, whether or not the quotients
are Hausdorff. In particular every finite-dimensional $W$ has closed
theta image in $\lthetaB$.
\end{lthetaproposition}

\begin{proof}
The range $R=\Theta V$ is closed in $\lthetaB$ and $\Theta:V\to R$ is a
homeomorphism. Therefore the closure of $\Theta W$ in $\lthetaB$ lies in
$R$ and equals its closure in $R$. A homeomorphism carries closures
to closures, proving \eqref{ltheta:eq:W}. Its restrictions to the closed
subspaces in that equality are inverse homeomorphisms. They carry
the equivalence classes modulo $W$ and $\Theta W$ to each other,
so the quotient universal property proves both induced maps
continuous and inverse.

For the last assertion, finite-dimensional subspaces of the present
Hausdorff locally convex space are closed. Here is a direct proof.
For a basis $w_1,\ldots,w_d$ of $W$, the map
$T:\lthetaC^d\to V$, $c\mapsto\sum c_iw_i$, is continuous and injective.
For each vector on the Euclidean unit sphere choose a continuous
seminorm positive on its image. Compactness gives finitely many
such seminorms whose sum $p$ is positive on every image on the
sphere; continuity and compactness give $p(Tc)\geq\delta|c|$ for
some $\delta>0$. Thus a convergent sequence $Tc_n$ in $V$ has
Cauchy coordinates, with limit $c$, and its limit is $Tc\in W$.
Metrizability gives closedness of $W$. Now apply \eqref{ltheta:eq:W}.
\end{proof}

Consequently for any programme subspace denoted $W_h$, formula
\eqref{ltheta:eq:W}, not an unproved replacement of $W_h$ by $V$, is the
exact available assertion. It computes its closure defect as the
source defect $\overline{W_h}^{\,V}/W_h$. For each fixed finite
polynomial source span $W_L$, its image is closed in $\lthetaB$ by the
finite-dimensional assertion. For a union of these spans the precise
formula remains
\[
 \overline{\bigcup_L\Theta W_L}^{\,\lthetaB}
    =\Theta\left(\overline{\bigcup_L W_L}^{\,V}\right)
\]
when the spans are increasing, so their union is a linear subspace.
No assertion that this source closure is all of $V$ is made.

The inclusion into the original Hilbert space
$\lthetaHH=L^2(\lthetaR_{>0},dx)$ is a different, explicitly continuous map:
\begin{equation}\label{ltheta:eq:Hilbert}
 \|F\|_{\lthetaHH}^2
 \leq p_{0,0}(F)^2\int_0^1dx
          +p_{1,0}(F)^2\int_1^\infty x^{-2}\,dx
 =p_{0,0}(F)^2+p_{1,0}(F)^2.
\end{equation}
It is injective since a continuous function zero almost everywhere
is zero everywhere. Its inverse on its image is not continuous.
Indeed fix $0\ne\eta\in C_c^\infty((1,2))$ and put
$F_n(x)=n^{-1}\eta(x-n)$. Each $F_n$ belongs to $\lthetaB$;
$\|F_n\|_{\lthetaHH}=n^{-1}\|\eta\|_2\to0$, whereas for a fixed
$t_0\in(1,2)$ with $\eta(t_0)\ne0$,
\[
 p_{2,0}(F_n)\geq n^{-1}(n+t_0)^2|\eta(t_0)|\longrightarrow\infty.
\]
Thus Hilbert convergence does not imply convergence in the original
target topology. In particular the closed-image theorem does not
assert Hilbert closedness, alter any finite Hilbert Gram matrix, or
replace a calculation of Hilbert closure. It removes exactly the
closure defect in (E5) for the original full source $V$ and the
original target seminorm topology.


\endgroup

\clearpage\begingroup
\part*{The actual arithmetic Jacobi commutant and its finite boundary}
\newtheorem{ljacobitheorem}{Theorem}
\newtheorem{ljacobiproposition}[ljacobitheorem]{Proposition}
\newcommand{\ljacobiR}{\mathbb R}
\newcommand{\ljacobiC}{\mathbb C}
\newcommand{\ljacobiN}{\mathbb N}
\newcommand{\ljacobiJ}{\mathcal J}
\newcommand{\ljacobiCC}{\mathcal C}
\newcommand{\ljacobiDom}{\operatorname{Dom}}



\section{Unchanged data and the proved envelope}

Retain the actual original packet divisor
\[
 \rho=\tfrac12+\delta+i\gamma,\quad 0<\delta<\tfrac12,\quad\gamma>2,
 \quad h(s)=\prod_{z\in\{\rho,\bar\rho,1-\rho,1-\bar\rho\}}(s-z)^m,
\]
with its full multiplicity $m\geq1$, $g=2\xi$, and
\[
 w_h(y)=\frac{|(g/h)(1/2+iy)|^2}{2\pi},\qquad
 m_k=w_h^{*k},\qquad \mu_h=\int_\ljacobiR w_h(y)\,dy.
\]
Here $k=4l+1\geq9$ is the original tensor order. These are the
original programme objects, and no existence assertion for an
off-line packet is added. The roots of $h$ are closed under conjugation
with unchanged multiplicities, so $h$ has real coefficients. The
original completed zeta function satisfies
$g(\bar s)=\overline{g(s)}$. Thus $w_h(-y)=w_h(y)$; convolution
preserves this equality and $m_k$ is even. Put
\[
 c=k/2,\quad q=[1+k(m-1)](k+1)^2,\qquad
 \chi(S)=\prod_{a,b=0}^k
 [S-c-(2a-k)\delta-i(2b-k)\gamma]^{1+k(m-1)}.
\]
The full finite quotient will use $N\geq q-1$.

The reference measure and its mass are exactly
\[
 d\sigma(y)=\rho_\sigma(y)\,dy
 =\frac{|\Gamma(1/4+iy/2)|^2}{2\pi}\,dy,
 \qquad M_\sigma=\sqrt{2\pi}.
\]
In particular, this measure is not replaced by a probability measure.
The already proved CAI16--18 envelope supplies
\begin{equation}\label{ljacobi:eq:envelope}
 a_k(1+y^2)^{-M}\rho_\sigma(y)\leq m_k(y)
                       \leq A_k\rho_\sigma(y),\qquad M=21+4m,
\end{equation}
where the retained constants are
\[
 \alpha=\pi/2,\quad p=8m-9/2,\quad B_h=42+8m,\quad
 \vartheta_h=\int_{-1}^1w_h(y)\,dy,
 \quad M_\alpha=\int_\ljacobiR e^{\alpha y}w_h(y)\,dy,
\]
\[
 c_\Gamma=\sqrt2e^{-7/6},\qquad C_\Gamma=\sqrt{2\sqrt5}e^{2/3},
\]
\[
 a_k=\frac{c_h\vartheta_h^{k-3}e^{-\alpha(k-3)}(k-2)^{-B_h}}
              {C_\Gamma2^{B_h/2}},\qquad
 A_k=\frac{C_h^{\rm tilt}}{c_\Gamma}k^{p+1}M_\alpha^{k-1}.
\]
The constants $c_h,C_h^{\rm tilt}$ are the original CAI lower and
upper envelope constants. Their proof is not replaced by a hypothesis:
the source is the retained programme provider
\texttt{CAI\_\allowbreak COMPLETE\_\allowbreak PREREQUISITE.tex},
equations CAI16--18.
The all-degree consequence, with
$\Delta_N=[1+4(N+M)^2]^M$, is
\begin{equation}\label{ljacobi:eq:finite-envelope}
 \frac{a_k}{\Delta_N}\|P\|_\sigma^2
       \leq\int|P(y)|^2m_k(y)\,dy\leq A_k\|P\|_\sigma^2
       \qquad(\deg P\leq N).
\end{equation}
The notation $\Delta_N$ here distinguishes the retained scalar CAI
constant from the source-coordinate matrix $D_N$ below.

The human literature input is Franti\v{s}ek \v{S}tampach and Pavel
\v{S}\v{t}ov\'{i}\v{c}ek, \emph{Spectral representation of some
weighted Hankel matrices and orthogonal polynomials from the Askey
scheme}; see
\href{https://arxiv.org/abs/1811.05820v1}{arXiv:1811.05820v1}.
We use Theorem \texttt{thm:commut\_\allowbreak matr\_\allowbreak Jac},
its following remark and proof,
and Theorem \texttt{thm:def\_H}. The following is a direct reconstruction
for the actual multiplier, including the mass absent from that
source's probability-measure convention.

\section{Gamma polynomials and the native Jacobi operator}

Define the monic polynomials by
\begin{equation}\label{ljacobi:eq:generating}
 G(y,t)=\sum_{n\geq0}p_n(y)\frac{t^n}{n!}
       =(1+t^2)^{-1/4}e^{y\arctan t}.
\end{equation}
Every branch here is its analytic branch at $t=0$. To verify the
mass and norms directly, the beta integral gives, for $a>0$,
\[
 \frac{\Gamma(a+i\tau)\Gamma(a-i\tau)}{\Gamma(2a)}
 =\int_\ljacobiR\frac{e^{i\tau v}}{(2\cosh(v/2))^{2a}}\,dv.
\]
This follows from $u=e^v$ in the beta integral on $(0,\infty)$.
Fourier inversion, followed by analytic continuation in the strip
where the Gamma envelope is integrable, gives at $a=1/4$ and $y=2\tau$
\begin{equation}\label{ljacobi:eq:mgf-sigma}
 \int_\ljacobiR e^{zy}\,d\sigma(y)
 =\sqrt{2\pi}(\cos z)^{-1/2},\qquad |\Re z|<\pi/2.
\end{equation}
In detail, inversion gives the Fourier transform at real frequency
$v$ as $2\Gamma(1/2)/(2\cosh v)^{1/2}$; setting $v=-iz$ yields
$\sqrt{2\pi}(\cos z)^{-1/2}$. The CAI17 Gamma envelope bounds the
integrand on every closed narrower strip and justifies holomorphy
and this continuation. At zero this proves the stated mass.

For sufficiently small $t,u$, the cosine addition formula and
\eqref{ljacobi:eq:mgf-sigma} imply
\[
 \int_\ljacobiR G(y,t)G(y,u)\,d\sigma(y)
 =M_\sigma(1-tu)^{-1/2}
 =M_\sigma\sum_{n\geq0}\frac{(1/2)_n}{n!}t^nu^n.
\]
All coefficient extractions may be taken on fixed small complex
circles, dominated by $Ce^{\varepsilon|y|}$ with $\varepsilon<\pi/2$.
Consequently
\[
 \int p_m p_n\,d\sigma=\delta_{mn}\gamma_n,
 \qquad\gamma_n=M_\sigma n!(1/2)_n.
\]
Define, without altering the measure,
\[
 \varphi_n=\frac{p_n}{\sqrt{\gamma_n}},\qquad
 r_n=\sqrt{M_\sigma}\varphi_n
     =\frac{p_n}{\sqrt{n!(1/2)_n}}.
\]
Thus $\varphi_0=M_\sigma^{-1/2}$, $r_0=1$, and
$\int r_mr_n\,d\sigma=M_\sigma\delta_{mn}$.
The $r_n$ are not orthonormal in $d\sigma$.

Differentiating \eqref{ljacobi:eq:generating} gives
$(1+t^2)\partial_tG=(y-t/2)G$ and hence
\[
 p_{n+1}=yp_n-n(n-1/2)p_{n-1},\qquad p_0=1,\quad p_1=y.
\]
Therefore
\begin{equation}\label{ljacobi:eq:recurrence}
 y\varphi_n=b_n\varphi_{n+1}+b_{n-1}\varphi_{n-1},\qquad
 b_n=\sqrt{(n+1)(n+1/2)},\quad b_{-1}=0,
\end{equation}
and the same recurrence holds for $r_n$. Also $p_n(-y)=(-1)^np_n(y)$.
The original source's Meixner--Pollaczek generating function gives
the exact dictionary
\[
 p_n(y)=n!P_n^{(1/4)}(y/2;\pi/2),\qquad
 r_n(y)=\sqrt{\frac{n!}{(1/2)_n}}P_n^{(1/4)}(y/2;\pi/2).
\]
Both the coordinate $y/2$ and every normalization factor remain.

The polynomials are complete in $L^2(\sigma)$. Indeed, if $f$ is
orthogonal to all polynomials, then $f\,d\sigma$ is a finite complex
measure, and Cauchy--Schwarz together with
\eqref{ljacobi:eq:mgf-sigma} shows that
$\int f(y)e^{zy}\,d\sigma(y)$ is holomorphic for $|\Re z|<\pi/4$.
Every derivative at zero is zero, so analytic uniqueness makes
the transform zero throughout that connected strip. Fourier
uniqueness on the imaginary axis makes $f\,d\sigma=0$, hence $f=0$.
The identical argument with the finite measure $(1+y^2)d\sigma$
proves polynomial density in that weighted space as well.

It follows that
\[
 U_\sigma:\ell^2(\ljacobiN_0)\longrightarrow L^2(\sigma),\qquad
 U_\sigma e_n=\varphi_n
\]
is unitary. The operator $J=U_\sigma^*M_yU_\sigma$ is self-adjoint,
and its matrix $\ljacobiJ$ has the entries \eqref{ljacobi:eq:recurrence}, with zero
diagonal. Polynomial density for $(1+y^2)d\sigma$ means that the
polynomials are a core of $M_y$ in its graph norm. Thus the finite
vectors are a core of $J$, so this is the unique self-adjoint closure
of the stated Jacobi matrix on finite vectors.

\section{The original bounded multiplier and exact first column}

Set $R_k=m_k/\rho_\sigma$. By \eqref{ljacobi:eq:envelope},
\[
 0<a_k(1+y^2)^{-M}\leq R_k(y)\leq A_k.
\]
Therefore multiplication by this exact density ratio is bounded,
positive and self-adjoint on $L^2(\sigma)$, and
\begin{equation}\label{ljacobi:eq:operator}
 \ljacobiCC=U_\sigma^*M_{R_k}U_\sigma,\qquad
 0\leq\ljacobiCC\leq A_k I,\qquad
 \ljacobiCC_{mn}=\int_\ljacobiR\varphi_m(y)\varphi_n(y)m_k(y)\,dy.
\end{equation}
This is not merely a formally defined matrix. The bounded multiplier
preserves $\ljacobiDom M_y$, since
$\|yR_k f\|_2\leq A_k\|yf\|_2$. Consequently $\ljacobiCC$ preserves
$\ljacobiDom J$ and $J\ljacobiCC=\ljacobiCC J$ there. Entrywise the same statement is
\[
 b_{m-1}\ljacobiCC_{m-1,n}+b_m\ljacobiCC_{m+1,n}
 =b_{n-1}\ljacobiCC_{m,n-1}+b_n\ljacobiCC_{m,n+1}.
\]
Each side is a finite sum; it also follows at once by integrating
\eqref{ljacobi:eq:recurrence} against the other polynomial and $m_k$.

Put $\alpha_j=\ljacobiCC_{j0}$. The exact native-mass expression is
\begin{equation}\label{ljacobi:eq:alpha-moment}
 \alpha_j=\frac1{M_\sigma}\int r_j(y)m_k(y)\,dy
 =\frac{\int p_j(y)m_k(y)\,dy}
             {M_\sigma\sqrt{j!(1/2)_j}}.
\end{equation}
In particular $\alpha_0=\mu_h^k/M_\sigma$ by Tonelli.

\begin{ljacobitheorem}\label{ljacobi:thm:first}
For every $m,n\geq0$ the complete entry is the finite sum
\begin{equation}\label{ljacobi:eq:first}
 \ljacobiCC_{mn}=\sum_{j=|m-n|}^{m+n}\alpha_j[r_j(\ljacobiJ)]_{mn}.
\end{equation}
Here $r_j(\ljacobiJ)$ is evaluated in the full Jacobi matrix, not in a
premature finite compression.
\end{ljacobitheorem}
\begin{proof}
The recurrence and $r_0=1$ give $r_n(\ljacobiJ)e_0=e_n$ by induction.
Indeed, for $n=0$ the recurrence gives $r_1(\ljacobiJ)e_0=e_1$; if the
assertion holds at $n,n-1$, then
$b_n r_{n+1}(\ljacobiJ)e_0=\ljacobiJ e_n-b_{n-1}e_{n-1}=b_ne_{n+1}$.
Expand the degree-$m+n$ polynomial $r_mr_n$ in the $r_j$ basis.
Orthogonality shows that its coefficient at $r_j$ is
\[
 d_{mnj}=\frac1{M_\sigma}\int r_mr_nr_j\,d\sigma
       =\int\varphi_m r_j\varphi_n\,d\sigma
       =[r_j(\ljacobiJ)]_{mn}.
\]
The last equality follows either from the multiplication operator
or by applying \eqref{ljacobi:eq:recurrence} finitely many times. Its
integral is zero when $j>m+n$ by orthogonality. It is also zero
when $m+j<n$ or $n+j<m$, by applying orthogonality to $r_n$ or
$r_m$ respectively. Thus only $|m-n|\leq j\leq m+n$ occurs.
Integrate this polynomial identity against $m_k/M_\sigma$ and
use \eqref{ljacobi:eq:alpha-moment}. The result is \eqref{ljacobi:eq:first}.
\end{proof}

Evenness of the unchanged weight gives $\alpha_j=0$ for odd $j$ and
$\ljacobiCC_{mn}=0$ for odd $m+n$. The first column also reconstructs the
actual multiplier in $L^2(\sigma)$, with the exact mass factor:
\begin{equation}\label{ljacobi:eq:mass-caveat}
 U_\sigma\ljacobiCC e_0=\frac{R_k}{\sqrt{M_\sigma}},\qquad
 R_k=\sqrt{M_\sigma}U_\sigma\ljacobiCC e_0
       =\sum_{j\geq0}\alpha_jr_j\quad\text{in }L^2(\sigma).
\end{equation}
To prove the last convergence, expand the vector $\ljacobiCC e_0\in\ell^2$
in its coordinates $\alpha_j$ and apply the unitary $U_\sigma$.
In particular the multiplier itself is not $U_\sigma\ljacobiCC e_0$ in
the original non-probability convention. No pointwise or
operator-norm convergence of the displayed series is asserted.

Define $M_h(z)=\int e^{zy}w_h(y)\,dy$. CAI16 gives
$M_\alpha<\infty$. Since the weight is even,
$\int e^{\alpha|y|}w_h\leq2M_\alpha$. Tonelli and convolution give
\[
 \int e^{zy}m_k(y)\,dy=M_h(z)^k,
 \qquad |\Re z|<\alpha.
\]
For real $z$ this is the absolutely convergent product integral;
the same domination gives the complex version and holomorphy.
For $|t|\leq1/4$,
$|\arctan t|\leq\operatorname{arctanh}(1/4)<\alpha$.
We can therefore integrate \eqref{ljacobi:eq:generating} on any smaller
complex circle and extract coefficients. This proves exactly
\begin{equation}\label{ljacobi:eq:alpha-generating}
 \boxed{\alpha_j=
 \frac{j!}{M_\sigma\sqrt{j!(1/2)_j}}
 [t^j]\left((1+t^2)^{-1/4}M_h(\arctan t)^k\right).}
\end{equation}
The factor $1/M_\sigma=1/\sqrt{2\pi}$ is necessary.

\section{The finite-cutoff sign and its full boundary vector}

Let $P_N$ project to $e_0,\ldots,e_N$, and write
\[
 C_N=P_N\ljacobiCC P_N,\quad J_N=P_N\ljacobiJ P_N,\quad
 c_N=(\ljacobiCC_{0,N+1},\ldots,\ljacobiCC_{N,N+1})^T.
\]
All are now finite coordinate matrices/vectors. Compressing the
full entrywise commutation identity gives
\[
 0=[J_N,C_N]+P_N\ljacobiJ(I-P_N)\ljacobiCC P_N
                    -P_N\ljacobiCC(I-P_N)\ljacobiJ P_N.
\]
The unique crossing Jacobi edge is $b_N$. Since $\ljacobiCC$ is real
symmetric, the two extra terms are respectively $b_Ne_Nc_N^*$
and $b_Nc_Ne_N^*$. Hence
\begin{equation}\label{ljacobi:eq:cutoff}
 \boxed{[J_N,C_N]=b_N(c_Ne_N^*-e_Nc_N^*).}
\end{equation}
This proves the stated sign. Formula \eqref{ljacobi:eq:first} computes
$C_N$ using the coefficients through index $2N$. It computes $c_N$
using those through index $2N+1$. Because the native odd coefficients
vanish, $\alpha_{2N+1}$ is zero; retaining it in the index range is
correct and harmless.

In fact evenness gives $(c_N)_N=\ljacobiCC_{N,N+1}=0$, so $c_N\perp e_N$.
If $c_N=0$, the finite matrices commute. If $c_N\ne0$, the
commutator vanishes on the orthogonal complement of
$\operatorname{span}\{e_N,c_N\}$ and acts on the ordered orthonormal
basis $e_N,c_N/\|c_N\|$ there as
\[
 \begin{pmatrix}0&-b_N\|c_N\|\\b_N\|c_N\|&0\end{pmatrix}.
\]
Thus its rank is two and its norm is exactly $b_N\|c_N\|$.
These are identities, not new estimates on the vector $c_N$.

The distinction between $P_Nr_j(\ljacobiJ)P_N$ and $r_j(J_N)$ is essential.
For example $r_2(y)=(y^2-1/2)/\sqrt{3/2}$, and at $N=1$,
\[
 [r_2(\ljacobiJ)]_{11}=\sqrt6,\qquad [r_2(J_1)]_{11}=0.
\]
Indeed $(\ljacobiJ^2)_{11}=b_0^2+b_1^2=1/2+3$, whereas
$(J_1^2)_{11}=b_0^2=1/2$. The missing term is the path from $1$
to $2$ and back to $1$. Accordingly
$\ljacobiCC_{11}=\alpha_0+\sqrt6\alpha_2$, not merely $\alpha_0$.

\section{Exact original source and observation coordinates}

For $F(S)=\sum_{j=0}^N a_jS^j$, the monomial coefficient vector
of the same function $F(c+iy)$ is $Ta$, where
\[
 T_{rj}=\binom jr c^{j-r}i^r\ (r\leq j),\qquad T_{rj}=0\ (r>j).
\]
This is the binomial theorem with no conjugation or mass change.
Let $B$ have the monomial coefficient vector of $\varphi_n$ as
column $n$. It is upper triangular with diagonal
$1/\sqrt{\gamma_n}$, while $T$ has diagonal $i^n$.
Both matrices are invertible. If $x$ denotes the Gamma-frame
coefficient vector of $F(c+iy)$, equality of the two polynomial
expansions is $Bx=Ta$. Hence
\begin{equation}\label{ljacobi:eq:source}
 x=D_Na,\qquad D_N=B^{-1}T,\qquad
 (D_N)_{nn}=i^n\sqrt{\gamma_n}.
\end{equation}
Integration in the unchanged arithmetic measure gives
\begin{equation}\label{ljacobi:eq:gram}
 H_{S,N}^{\rm ar}=D_N^*C_ND_N.
\end{equation}
For every $x\in\ljacobiC^{N+1}$, \eqref{ljacobi:eq:finite-envelope} applied to
$\sum x_n\varphi_n$ gives the concrete positive bounds
\begin{equation}\label{ljacobi:eq:CN-bounds}
 \frac{a_k}{\Delta_N}I\leq C_N\leq A_kI.
\end{equation}
Thus the inverse in the next formula exists for every fixed $N$.
No positive lower bound independent of $N$ for the infinite
operator is inferred from the decaying lower density envelope.

Retain the actual observation map
$\mathcal A_S=\Lambda\mathcal J_{S,N}$, including its full physical
Taylor unit, in its original source and observation coordinates.
Its codomain is the original observation image with its retained
Hermitian structure, so it is onto. Its matrix in the Gamma-frame
source coordinates must be
\begin{equation}\label{ljacobi:eq:E}
 \mathcal E_N=\mathcal A_SD_N^{-1},
\end{equation}
since $\mathcal A_Sa=\mathcal A_SD_N^{-1}x$. The symbol
$\mathcal J_{S,N}$ denotes the original class map, not the Jacobi
matrix $J_N$.

\begin{ljacobitheorem}\label{ljacobi:thm:quotient}
For every observed vector $u$ the original attained minimum is
\begin{equation}\label{ljacobi:eq:quotient}
 \min_{\mathcal A_Sa=u}a^*H_{S,N}^{\rm ar}a
  =u^*Q_{B,N}u,\qquad
 Q_{B,N}=(\mathcal E_NC_N^{-1}\mathcal E_N^*)^{-1},
\end{equation}
and its unique minimizing source vector is
\begin{equation}\label{ljacobi:eq:lift}
 a_{\min}=D_N^{-1}C_N^{-1}\mathcal E_N^*Q_{B,N}u.
\end{equation}
\end{ljacobitheorem}
\begin{proof}
Surjectivity of $\mathcal E_N$ makes its adjoint injective, so
$K=\mathcal E_NC_N^{-1}\mathcal E_N^*$ is strictly positive on the
observation space. Put $x_0=C_N^{-1}\mathcal E_N^*K^{-1}u$.
Then $\mathcal E_Nx_0=u$ and
$x_0^*C_Nx_0=u^*K^{-1}u$. Every other lift is $x_0+z$ with
$\mathcal E_Nz=0$, and
$z^*C_Nx_0=(\mathcal E_Nz)^*K^{-1}u=0$. Therefore its energy is
$u^*K^{-1}u+z^*C_Nz$, strictly larger unless $z=0$.
The bijection $x=D_Na$ and \eqref{ljacobi:eq:gram} prove both formulas in
the original source coordinates. Equivalently,
\[
 H_{S,N}^{{\rm ar},-1}=D_N^{-1}C_N^{-1}(D_N^*)^{-1},\qquad
 \mathcal A_SH_{S,N}^{{\rm ar},-1}\mathcal A_S^*
                   =\mathcal E_NC_N^{-1}\mathcal E_N^*.
\]
Thus this is exactly the original quotient metric, not an auxiliary
minimum for a changed source.
\end{proof}

For completeness, \eqref{ljacobi:eq:CN-bounds} also gives the exact matrix
comparison
\[
 \frac{a_k}{\Delta_N}(\mathcal E_N\mathcal E_N^*)^{-1}
       \leq Q_{B,N}\leq
 A_k(\mathcal E_N\mathcal E_N^*)^{-1}.
\]
This follows by inverting the two positive bounds for $C_N$,
congruence with $\mathcal E_N$, and inversion once again. It does
not replace $\mathcal E_N$ by another observation map.

\section{Explicit finite checks and retained phase problem}

Write $\nu_j=\int y^jm_k(y)\,dy$ and
$\eta_j=\int y^jw_h(y)\,dy$, so $\eta_0=\mu_h$.
The first polynomial values from the exact recurrence are
\[
 p_0=1,\quad p_1=y,\quad p_2=y^2-\tfrac12,\quad
 p_3=y^3-\tfrac72y,\quad p_4=y^4-11y^2+\tfrac{15}4.
\]
Consequently the first even column entries are
\[
 \alpha_0=\frac{\nu_0}{M_\sigma},\qquad
 \alpha_2=\frac{\nu_2-\nu_0/2}{M_\sigma\sqrt{3/2}},\qquad
 \alpha_4=\frac{\nu_4-11\nu_2+15\nu_0/4}
                   {M_\sigma\sqrt{315/2}}.
\]
Multinomial expansion of the actual convolution, using evenness,
gives
\[
 \nu_0=\mu_h^k,\quad \nu_2=k\eta_2\mu_h^{k-1},\quad
 \nu_4=k\eta_4\mu_h^{k-1}
              +3k(k-1)\eta_2^2\mu_h^{k-2}.
\]
The $3k(k-1)$ factor counts the two distinct factors contributing
second powers: $\binom{k}{2}\binom42=3k(k-1)$.

At cutoff $N=1$, evenness gives
\[
 C_1=\frac1{M_\sigma}\begin{pmatrix}\nu_0&0\\0&2\nu_2\end{pmatrix},
 \quad J_1=\begin{pmatrix}0&1/\sqrt2\\1/\sqrt2&0\end{pmatrix},
 \quad c_1=\begin{pmatrix}\alpha_2\\0\end{pmatrix},\quad b_1=\sqrt3.
\]
Its $(0,1)$ commutator entry is
$(2\nu_2-\nu_0)/(\sqrt2M_\sigma)=\sqrt3\alpha_2$,
confirming both the sign and the mass factor in \eqref{ljacobi:eq:cutoff}.
The degree-two source change is explicitly
\[
 D_2=\sqrt{M_\sigma}
 \begin{pmatrix}
 1&c&c^2-1/2\\
 0&i/\sqrt2&\sqrt2ic\\
 0&0&-\sqrt{3/2}
 \end{pmatrix}.
\]
In particular the source constant $F=1$ has coefficient
$x_0=\sqrt{M_\sigma}$ and energy
$M_\sigma(C_0)_{00}=\mu_h^k$, retaining the original arithmetic mass.

The companion exact symbolic script \texttt{check\_symbolic.py}
checks \eqref{ljacobi:eq:first} for all $0\leq m,n\leq4$ (25 entries),
checks \eqref{ljacobi:eq:cutoff} at $N=0,1,2,3$, checks the generating
polynomials through degree six, and checks
\eqref{ljacobi:eq:source}--\eqref{ljacobi:eq:gram} through degree two. All moment
variables are independent real symbols; the checks even allow
nonzero odd moments. The parameter representing $M_\sigma$ remains
a symbolic positive mass rather than being set to one. These
finite checks supplement, and do not replace, the general proofs.

For the original classes $a_1=[p_N^{\rm ar}]_\chi$,
$a_2=[p_{N+1}^{\rm ar}]_\chi$ and original terminal class $v$, the
retained quantities are therefore exactly
\[
 B_i=(\Lambda a_i)^*Q_{B,N}\Lambda v,\qquad
 \Phi_{B,N}=\frac{2\Re(\overline{B_1}B_2)}{\omega_{N,k}^{\rm ar}}.
\]
The formulas above give a finite representation of the same
quantities and the full Jacobi compression defect. They do not
evaluate or estimate these complex signed products. The physical
Taylor unit, the observation map, and both original classes remain
present; parity of $C_N$ does not remove them.


\endgroup

% END INDEXED LITERATURE EXACT RECEIVERS 20260920


% BEGIN COMPLETE PARITY AND CONDUCTOR RECEIVERS

\clearpage\begingroup
\part*{Original parity and Pollaczek source bounds}
\newtheorem{paritytheorem}{Theorem}
\newtheorem{parityproposition}[paritytheorem]{Proposition}
\newtheorem{paritylemma}[paritytheorem]{Lemma}
\newcommand{\parityR}{\mathbb R}
\newcommand{\parityC}{\mathbb C}
\newcommand{\parityN}{\mathbb N}
\newcommand{\parityPol}{\mathbb C[x]}
\newcommand{\paritydiag}{\operatorname{diag}}
\newcommand{\parityRan}{\operatorname{Ran}}



\section{The unchanged arithmetic weight and the literature input}

Retain the original quartet, divisor, and completed numerator
\[
 \rho=\tfrac12+\delta+i\gamma,\quad 0<\delta<\tfrac12,\quad\gamma>2,
 \quad h(s)=\prod_{z\in\{\rho,\bar\rho,1-\rho,1-\bar\rho\}}(s-z)^m,
 \quad m\geq1,\quad g=2\xi,
\]
\[
 w_h(y)=\frac{|(g/h)(1/2+iy)|^2}{2\pi},\quad
 m_k=w_h^{*k},\quad k=4l+1\geq9,\quad
 \mu_h=\int_\parityR w_h(y)\,dy.
\]
All multiplicities are the original complete multiplicities. No
existence of an off-line quartet is asserted here. Conjugation
symmetry of $g$ and of the root set of $h$ gives $w_h(-y)=w_h(y)$,
and convolution gives the same symmetry for $m_k$.
The source variable, centre and quotient polynomial remain
\[
 S=c+iy,\quad c=k/2,\quad q=[1+k(m-1)](k+1)^2,
\]
\[
 \chi(S)=\prod_{a,b=0}^k
 [S-c-(2a-k)\delta-i(2b-k)\gamma]^{1+k(m-1)}.
\]
All finite observation formulas below use the original degree
$N\geq q-1$ and actual map $\mathcal A_S=\Lambda\mathcal J_{S,N}$,
with its full physical Taylor unit. Its codomain is its original
observation image with the retained Hermitian structure.

The original proved CAI16--18 envelope is used at its actual scope
\cite{parity:CAI}; the arbitrary-fixed-multiplicity form is HG21 \cite{parity:HG}:
\begin{equation}\label{parity:eq:CAI}
 a_k(1+y^2)^{-M}\rho_\sigma(y)\leq m_k(y)\leq A_k\rho_\sigma(y),
 \quad M=21+4m,
\end{equation}
\[
 \rho_\sigma(y)=\frac{|\Gamma(1/4+iy/2)|^2}{2\pi},\quad
 \int\rho_\sigma=\sqrt{2\pi},\quad\alpha=\pi/2,
\]
\[
 c_\Gamma e^{-\alpha|y|}(1+|y|)^{-1/2}
 \leq\rho_\sigma(y)\leq
 C_\Gamma e^{-\alpha|y|}(1+|y|)^{-1/2},
\]
\[
 c_\Gamma=\sqrt2e^{-7/6},\quad C_\Gamma=\sqrt{2\sqrt5}e^{2/3},
 \quad p=8m-9/2,\quad B_h=42+8m,
\]
\[
 a_k=\frac{c_h\vartheta_h^{k-3}e^{-\alpha(k-3)}(k-2)^{-B_h}}
                   {C_\Gamma2^{B_h/2}},\quad
 A_k=\frac{C_h^{\rm tilt}}{c_\Gamma}k^{p+1}M_\alpha^{k-1},
\]
\[
 \vartheta_h=\int_{-1}^1w_h(y)\,dy,\qquad
 M_\alpha=\int_\parityR e^{\alpha y}w_h(y)\,dy.
\]
The constants $c_h,C_h^{\rm tilt}$ are precisely the established
lower-envelope and compact/tail constants in the complete CAI
provider, not new assumptions. In particular $m_k>0$ and all
moments exist. For complex polynomials $P$ of degree at most $d$,
that provider proves
\begin{equation}\label{parity:eq:CAI-finite}
 \frac{a_k}{\Delta_d}\|P\|_\sigma^2
       \leq\|P\|_{m_k}^2\leq A_k\|P\|_\sigma^2,
 \qquad \Delta_d=[1+4(d+M)^2]^M.
\end{equation}
The Gamma mass has not been divided out.

The original author source used here is A.~I.~Aptekarev,
G.~L\'opez Lagomasino and A.~Mart\'{\i}nez--Finkelshtein,
\emph{Strong asymptotics for the Pollaczek multiple orthogonal
polynomials ensembles},
\href{https://arxiv.org/abs/1410.1261v1}{arXiv:1410.1261v1}
\cite{parity:ALM}.
Its Introduction gives the weights
\begin{equation}\label{parity:eq:literature}
 w_1(x)=\frac1{\sinh(\alpha\sqrt x)},\qquad
 w_2(x)=\frac1{\sqrt x\cosh(\alpha\sqrt x)},\qquad x>0,
\end{equation}
and their positive Stieltjes ratio
\begin{equation}\label{parity:eq:S}
 s(x)=\frac{w_2(x)}{w_1(x)}
     =\frac{\tanh(\alpha\sqrt x)}{\sqrt x}
     =\frac4\pi\sum_{j=0}^{\infty}\frac1{x+(2j+1)^2}.
\end{equation}
The lowercase $s(x)$ in this note denotes this multiplier, not
the original complex source variable $S$. No multiple-orthogonal
asymptotic is assumed for the native weight.

\section{Exact parity-square maps, source images, and masses}

Define the two original transported measures on $(0,\infty)$ by
\begin{equation}\label{parity:eq:nu}
 d\nu_0(x)=\frac{m_k(\sqrt x)}{\sqrt x}\,dx,
 \qquad d\nu_1(x)=\sqrt x\,m_k(\sqrt x)\,dx=x\,d\nu_0(x).
\end{equation}
They have strictly positive densities. Their literal density
ratios against the first literature weight are
\begin{equation}\label{parity:eq:ratios}
 r_0(x)=\frac{m_k(\sqrt x)\sinh(\alpha\sqrt x)}{\sqrt x},
 \qquad r_1(x)=x r_0(x),\qquad d\nu_j=r_jw_1\,dx.
\end{equation}
For every polynomial $f(y)$ there is a unique pair $E,O$ with
$f(y)=E(y^2)+yO(y^2)$. Evenness of $m_k$ makes the cross term in
the squared norm odd. The substitution $x=y^2$ on each half-axis
then proves the exact sesquilinear and norm identities
\begin{equation}\label{parity:eq:parity-norm}
 \|f\|_{m_k}^2=\int_0^\infty|E(x)|^2\,d\nu_0(x)
                    +\int_0^\infty|O(x)|^2\,d\nu_1(x).
\end{equation}
There is no missing factor two: the two half-axes cancel the
$1/2$ in $dy=dx/(2\sqrt x)$.

The same construction is a unitary isomorphism
\[
 L^2(\parityR,m_kdy)\longrightarrow L^2(\nu_0)\oplus L^2(\nu_1),
\]
with
\[
 E(x)=\frac{f(\sqrt x)+f(-\sqrt x)}2,\qquad
 O(x)=\frac{f(\sqrt x)-f(-\sqrt x)}{2\sqrt x}.
\]
These formulas are defined almost everywhere, and their inverse
is $f(y)=E(y^2)+yO(y^2)$ for $y\ne0$. The identity of norms proves
both that the maps are well defined and that they are inverse
isometries.

The exact isometry into the literature's two Hilbert spaces is
\begin{equation}\label{parity:eq:literature-isometry}
 (E,O)\longmapsto
       \left(\sqrt{r_0}\,E,\sqrt{r_1/s}\,O\right)
 \quad\text{in }L^2(w_1dx)\oplus L^2(w_2dx).
\end{equation}
Since all density ratios are positive, division by their square
roots gives its inverse on these full Hilbert spaces. On the
degree-$N$ source its image is exactly
\[
 \sqrt{r_0}\,\parityPol_{\lfloor N/2\rfloor}
       \ \oplus\ \sqrt{r_1/s}\,\parityPol_{\lfloor(N-1)/2\rfloor},
\]
not an unweighted polynomial pair. Here $\parityPol_n$ means degree at
most $n$, and $\parityPol_{-1}=\{0\}$. The common-first-weight version
is $(E,O)\mapsto(\sqrt{r_0}E,\sqrt{r_1}O)$ in two copies of
$L^2(w_1dx)$.

Write $\eta_j=\int y^jw_h(y)\,dy$ and
$\zeta_j=\int y^jm_k(y)\,dy$. Tonelli and the convolution
multinomial formula give the original parity masses
\begin{equation}\label{parity:eq:masses}
 \nu_0((0,\infty))=\zeta_0=\mu_h^k,
 \qquad \nu_1((0,\infty))=\zeta_2=k\eta_2\mu_h^{k-1}.
\end{equation}
The cross terms in the second formula vanish because $w_h$ is
even. The literature masses are different fixed numbers:
\begin{equation}\label{parity:eq:raw-moments}
 \tau_j:=\int_0^\infty x^j w_1(x)\,dx
 =\frac{4(2j+1)!}{\alpha^{2j+2}}
        (1-2^{-2j-2})\zeta(2j+2),\qquad \tau_0=2,
\end{equation}
\[
 \int_0^\infty w_2(x)\,dx=2.
\]
For the first formula use $x=y^2$ and the positive expansion
$1/\sinh(\alpha y)=2\sum_{j\geq0}e^{-(2j+1)\alpha y}$,
then integrate each power of $y$. Its constant case uses
$\sum_{j\geq0}(2j+1)^{-2}=\pi^2/8$, also obtained below from
\eqref{parity:eq:S} at zero. For the second mass, the same substitution
gives $2\int_0^\infty\operatorname{sech}(\alpha y)dy=\pi/\alpha=2$;
an antiderivative of $\operatorname{sech} u$ is
$\arctan(\sinh u)$. None of these masses is set equal to one.

For the exact original source $F(S)=\sum_{j=0}^Na_jS^j$, define
\[
 T_{rj}=\binom jr c^{j-r}i^r\quad(r\leq j),\qquad T_{rj}=0\quad(r>j).
\]
Let $\Pi_N$ reorder the coefficients in powers of $y$ into even
powers first and odd powers second, without rescaling any entry.
The parity coefficient map is the invertible matrix
\begin{equation}\label{parity:eq:source-map}
 L_N=\Pi_NT,\qquad (e,o)=L_Na,
\end{equation}
where explicitly
\[
 e_r=(-1)^r\sum_{j=2r}^N\binom j{2r}c^{j-2r}a_j,
 \quad
 o_r=i(-1)^r\sum_{j=2r+1}^N\binom j{2r+1}c^{j-2r-1}a_j.
\]
Thus the full original degree-$N$ source maps onto precisely the
displayed pair of polynomial spaces. An original constrained
source subspace $W$ maps to $L_NW$, not to a substituted full pair.

Put $n_0=\lfloor N/2\rfloor$, $n_1=\lfloor(N-1)/2\rfloor$ and
\begin{equation}\label{parity:eq:parity-grams}
 H_{j,n}=\left[\int_0^\infty x^{a+b}\,d\nu_j(x)\right]_{a,b=0}^n,
 \qquad H=\paritydiag(H_{0,n_0},H_{1,n_1}).
\end{equation}
The matrices are positive definite because nonzero polynomials
cannot vanish on a set of full measure for a positive density.
They retain the literal moments
$(H_{j,n})_{ab}=\zeta_{2(a+b+j)}$. Integration gives
\begin{equation}\label{parity:eq:original-gram}
 H_{S,N}^{\rm ar}=L_N^*HL_N.
\end{equation}
The actual observation matrix in these coordinates is
\begin{equation}\label{parity:eq:observation}
A=\mathcal A_SL_N^{-1}.
\end{equation}
The exact original observation metric and minimizing lift are
\begin{equation}\label{parity:eq:Q}
 Q(H)=(AH^{-1}A^*)^{-1},\qquad
 a_{\min}=L_N^{-1}H^{-1}A^*Q(H)u.
\end{equation}
Indeed $A$ is onto its stated observation image, so $AH^{-1}A^*$
is strictly positive. The vector $z_0=H^{-1}A^*Q(H)u$ satisfies
$Az_0=u$, and for $Az=0$ the cross term is
$z^*Hz_0=(Az)^*Q(H)u=0$. Every other lift has energy
$u^*Q(H)u+z^*Hz$, proving the formula and uniqueness. The
congruence \eqref{parity:eq:original-gram} returns the result to the
unchanged source coefficients.

There is also an exact transport of the original quotient, not
only of its source norm. Put
\[
 u_a=(2a-k)\delta,\qquad v_b=(2b-k)\gamma,\qquad
 e_k=1+k(m-1),
\]
and define the monic polynomial
\begin{equation}\label{parity:eq:Xi}
 \Xi(x)=\prod_{a=0}^k\prod_{b=(k+1)/2}^{k}
                 [x-(v_b-iu_a)^2]^{e_k}.
\end{equation}
Its degree is $q/2$. The opposite grid point of $(a,b)$ is
$(k-a,k-b)$, and exactly one point of each opposite pair has
$v_b>0$. That pair's factors after $S=c+iy$ are
\[
 (iy-u_a-iv_b)(iy+u_a+iv_b)
          =-[y^2-(v_b-iu_a)^2].
\]
Since $k$ is odd, $q=e_k(k+1)^2$ is divisible by four. The number
$q/2$ of paired factors with multiplicity is therefore even,
so the full sign is $(-1)^{q/2}=1$. Consequently
\begin{equation}\label{parity:eq:chi-square}
 \chi(c+iy)=\Xi(y^2)
\end{equation}
with exactly the original leading coefficient. Conjugating the
roots in \eqref{parity:eq:Xi} is the permutation $a\mapsto k-a$, so
$\Xi$ also has real coefficients.

Let $\mathcal E_x=\parityC[x]/(\Xi)$. Substitution and parity give the
explicit algebra isomorphism followed by vector-space coordinates
\begin{equation}\label{parity:eq:quotient-parity}
 \parityC[S]/(\chi)\longrightarrow\parityC[y]/(\Xi(y^2))
                  \longrightarrow\mathcal E_x\oplus\mathcal E_x,
 \qquad [F]\longmapsto([E],[O]).
\end{equation}
The second arrow has multiplication
$(E,O)(E',O')=(EE'+xOO',EO'+OE')$ modulo $\Xi$ in each component.
To verify its kernel, divide $E$ and $O$ separately by $\Xi$.
The remaining even and odd polynomials have degrees less than
$q$ in $y$, so their sum is divisible by the monic degree-$q$
polynomial $\Xi(y^2)$ only when both remainders are zero. This
also proves surjectivity and identifies the inverse through the
unique remainder $E(y^2)+yO(y^2)$ and $y=(S-c)/i$.
Multiplication by the original variable $S$ is precisely
\begin{equation}\label{parity:eq:S-block}
 \begin{pmatrix}cI&iM_x\\ iI&cI\end{pmatrix}
 \quad\text{on }\mathcal E_x\oplus\mathcal E_x.
\end{equation}
An original full Taylor unit $\mathfrak u\in\parityC[S]/(\chi)$,
transported as $(U_0,U_1)$, acts by the unchanged block matrix
\begin{equation}\label{parity:eq:unit-block}
 \begin{pmatrix}M_{U_0}&M_{xU_1}\\M_{U_1}&M_{U_0}\end{pmatrix}.
\end{equation}
Its inverse is the transported multiplication by
$\mathfrak u^{-1}$; the unit has not been removed. Thus the
original class and observation maps can equivalently be
transported through \eqref{parity:eq:quotient-parity} and
\eqref{parity:eq:unit-block}. Formula \eqref{parity:eq:observation} is their
same matrix composition before choosing these quotient
coordinates. At $N\geq q-1$, both parity source degree bounds
are at least $q/2-1$, so both residue-class components are
attained by the original degree-$N$ source.

\section{Both native orthogonal families and all Christoffel constants}

Let $P_d(y)$ be the real monic orthogonal polynomial of degree $d$
for $m_k(y)dy$ and let $\omega_d=\int P_d^2m_k>0$. Existence and
uniqueness follow by solving the positive definite moment system.
Evenness and uniqueness give $P_d(-y)=(-1)^dP_d(y)$, so define
\[
 P_{2n}(y)=E_n(y^2),\quad P_{2n+1}(y)=yO_n(y^2),
 \quad \kappa_n=\omega_{2n},\quad\lambda_n=\omega_{2n+1}.
\]
Equation \eqref{parity:eq:parity-norm} shows that $E_n$ and $O_n$ are
respectively the real monic orthogonal families for $\nu_0$ and
$\nu_1$, with norms $\kappa_n$ and $\lambda_n$.

The monic orthogonal polynomial in the original complex source
coordinate is not obtained by dropping the phase. It is exactly
\begin{equation}\label{parity:eq:S-polynomial-phase}
 \mathsf P_d(S)=i^dP_d((S-c)/i),\qquad
 \mathsf P_d(c+iy)=i^dP_d(y),\qquad
 \|\mathsf P_d\|_{m_k}^2=\omega_d.
\end{equation}
Its leading coefficient is $i^di^{-d}=1$; the lower-degree
orthogonality follows by the same invertible substitution used
in \eqref{parity:eq:source-map}. Thus uniqueness proves this exact
dictionary. Its parity pair is $((-1)^nE_n,0)$ at degree $2n$
and $(0,i(-1)^nO_n)$ at degree $2n+1$. These phases remain in
the original boundary classes and observation entries.

Orthogonality gives the full original recurrence
\[
 yP_d=P_{d+1}+\beta_dP_{d-1},\qquad
 \beta_d=\omega_d/\omega_{d-1}>0\quad(d\geq1),\quad \beta_0=0.
\]
To prove it, expand $yP_d-P_{d+1}$ in the lower orthogonal basis.
Its coefficients below $d-1$ vanish by transferring $y$ to the
lower polynomial. The coefficient of $P_d$ vanishes by parity,
and the coefficient of $P_{d-1}$ is
$\langle P_d,yP_{d-1}\rangle/\omega_{d-1}=\omega_d/\omega_{d-1}$.
In the original $S$ coordinate the same identity reads
$(S-c)\mathsf P_d=\mathsf P_{d+1}-\beta_d\mathsf P_{d-1}$,
since $i^{d+1}/i^{d-1}=-1$; its sign is therefore retained too.

Set
\begin{equation}\label{parity:eq:ad}
 a_n=\frac{\lambda_n}{\kappa_n}=\beta_{2n+1}>0,
 \qquad d_n=\frac{\kappa_n}{\lambda_{n-1}}=\beta_{2n}>0\ (n\geq1),
 \qquad d_0=0.
\end{equation}
Separating the even and odd instances of that recurrence gives
the exact polynomial identities
\begin{equation}\label{parity:eq:Christoffel}
 xO_n=E_{n+1}+a_nE_n,\qquad
 E_n=O_n+d_nO_{n-1},\quad O_{-1}=0.
\end{equation}
Evaluating the first at zero proves, starting from $E_0=1$,
\begin{equation}\label{parity:eq:E0}
 E_n(0)=(-1)^n\prod_{j=0}^{n-1}a_j\ne0,
 \qquad \frac{E_{n+1}(0)}{E_n(0)}=-a_n.
\end{equation}
Therefore the exact Christoffel formula is
\begin{equation}\label{parity:eq:Christoffel-explicit}
 O_n(x)=\frac{E_{n+1}(x)-[E_{n+1}(0)/E_n(0)]E_n(x)}x
        =\frac{E_{n+1}(x)+a_nE_n(x)}x,
 \quad \lambda_n=a_n\kappa_n.
\end{equation}
The numerator has a zero at zero, so this is a polynomial identity,
including the endpoint. The complete odd zero value is
\begin{equation}\label{parity:eq:O0}
 O_n(0)=(-1)^n\sum_{j=0}^n
     \left(\prod_{r=j+1}^nd_r\right)
     \left(\prod_{r=0}^{j-1}a_r\right)
 =E_{n+1}'(0)+a_nE_n'(0).
\end{equation}
The sum follows by iterating $O_n(0)=E_n(0)-d_nO_{n-1}(0)$;
the derivative formula follows from \eqref{parity:eq:Christoffel-explicit}.
Every empty product equals one, so this covers $n=0$.

Substitution of the two identities in \eqref{parity:eq:Christoffel} into
each other yields both complete monic recurrences
\begin{align}
 xE_n&=E_{n+1}+(a_n+d_n)E_n+a_{n-1}d_nE_{n-1},\label{parity:eq:E-rec}\\
 xO_n&=O_{n+1}+(a_n+d_{n+1})O_n+a_nd_nO_{n-1}.\label{parity:eq:O-rec}
\end{align}
The last term of the first is zero at $n=0$ and needs no $a_{-1}$.
The norm ratios are exactly
$\kappa_n/\kappa_{n-1}=a_{n-1}d_n$ and
$\lambda_n/\lambda_{n-1}=a_nd_n$.
For example the orthonormal version of the first cross relation is
\[
 x\frac{O_n}{\sqrt{\lambda_n}}
  =\sqrt{d_{n+1}}\frac{E_{n+1}}{\sqrt{\kappa_{n+1}}}
       +\sqrt{a_n}\frac{E_n}{\sqrt{\kappa_n}},
\]
and the second is
\[
 \frac{E_n}{\sqrt{\kappa_n}}
 =\sqrt{a_n}\frac{O_n}{\sqrt{\lambda_n}}
       +\sqrt{d_n}\frac{O_{n-1}}{\sqrt{\lambda_{n-1}}}.
\]
These equalities are between functions in the displayed different
weighted spaces; they explicitly give the matrix of the identity
map or $x$-multiplication on every finite polynomial domain.

\section{The exact multiple-orthogonality receiver and its residuals}

The raw paper has one polynomial $P$ subject to moments against
both $w_1$ and $w_2$. The native common-multiplier Nikishin pair is
\[
 (d\nu_0,\ s\,d\nu_0)=r_0(x)(w_1(x),w_2(x))\,dx;
\]
the native odd measure is instead $d\nu_1=x\,d\nu_0$.
The literal maps between their functionals, for every polynomial
$P$ and nonnegative integer $j$, are
\begin{align}
 \int x^jP\,d\nu_0&=\int x^j(r_0P)w_1\,dx,\label{parity:eq:functional0}\\
 \int x^jP\,s\,d\nu_0&=\int x^j(r_0P)w_2\,dx,\label{parity:eq:functionalS}\\
 \int x^jP\,d\nu_1&=\int x^j\left(\frac{r_1}{s}P\right)w_2\,dx
                    =\int x^{j+1}P\,d\nu_0.\label{parity:eq:functional1}
\end{align}
Thus each finite raw moment operator receives the explicit
multiplied source space in these formulas, not the raw polynomial
space unless that multiplier is retained. These maps are the
functional counterparts of the square-root isometry
\eqref{parity:eq:literature-isometry}.

Here is a nonzero residual computation within the native common
multiplier pair. The $n$ zeros of $E_n$ are simple and lie in
$(0,\infty)$. To prove this, multiply $E_n$ by the product of its
distinct sign-changing zeros in that interval. If fewer than $n$
such zeros existed, this product would have degree less than $n$
and its product with $E_n$ would have constant nonzero sign except
at finitely many points of the positive interval. Positivity of
$\nu_0$ would contradict orthogonality. Thus there are $n$ such
zeros, which exhaust the degree and establish the assertion.
In particular $E_n(-a)$ has sign $(-1)^n$ for every $a>0$.

For $a>0$, divide $E_n(x)-E_n(-a)$ by $x+a$. The quotient has
degree $n-1$, so orthogonality gives the exact positive-resolvent
identity
\begin{equation}\label{parity:eq:residual-atom}
 \int\frac{E_n(x)}{x+a}\,d\nu_0(x)
  =\frac1{E_n(-a)}\int\frac{E_n(x)^2}{x+a}\,d\nu_0(x).
\end{equation}
For $0\leq j<n$, division of $x^j$ by $x+a$ similarly gives
\begin{equation}\label{parity:eq:residual}
 \int x^jE_n(x)s(x)\,d\nu_0(x)
 =\frac4\pi\sum_{r=0}^{\infty}
 \frac{[-(2r+1)^2]^j}{E_n(-(2r+1)^2)}
 \int\frac{E_n(x)^2}{x+(2r+1)^2}\,d\nu_0(x).
\end{equation}
To justify the sum before division, dominate by
$\int|x^jE_n(x)|s(x)\,d\nu_0<\infty$; the atomwise integrals are
absolutely summable since $s$ is bounded. The divisions preserve
each integral exactly. Every term on the resulting right side
has the same strict sign $(-1)^{n+j}$, so that is the sign of the
nonzero residual. In particular $E_n$ does not satisfy even the
zeroth second-weight condition of the native common-multiplier
Nikishin system. This is an evaluated residual, not just a
declaration that two orthogonalities differ. For $n=0$ the
zeroth residual is directly positive.

There is also a completely finite lower bound. Cauchy--Schwarz
with the factors $(x+1)^{1/2}E_n$ and $(x+1)^{-1/2}E_n$ yields
\[
 \int\frac{E_n^2}{x+1}\,d\nu_0
 \geq\frac{\kappa_n^2}{\int(x+1)E_n^2\,d\nu_0}
 =\frac{\kappa_n}{1+a_n+d_n},
\]
where \eqref{parity:eq:E-rec} computes the final moment. The first
positive atom in \eqref{parity:eq:residual} therefore gives
\begin{equation}\label{parity:eq:residual-lower}
 \left|\int E_n s\,d\nu_0\right|
 \geq\frac{4\kappa_n}{\pi(1+a_n+d_n)|E_n(-1)|}>0.
\end{equation}
All its constants belong to the original native parity family.

\section{Two positive resolvent series, with all constants proved}

We give a direct proof of \eqref{parity:eq:S} and the reciprocal identity
that will supply finite native bounds. On the interval
$[0,\alpha]$, $\alpha=\pi/2$, the operator $-d^2/dt^2+x$ with
Dirichlet condition at zero and Neumann condition at $\alpha$
has eigenvalues $x+(2j+1)^2$ and normalized eigenfunctions
$\sqrt{2/\alpha}\sin((2j+1)t)$. Their completeness is the ordinary
sine Fourier expansion obtained by odd reflection at zero and
even reflection at $\alpha$. The Green kernel is
\[
 G_x(u,v)=\frac{\sinh(\sqrt x\min(u,v))
                  \cosh(\sqrt x(\alpha-\max(u,v)))}
                    {\sqrt x\cosh(\alpha\sqrt x)}.
\]
Its formula follows by solving the homogeneous differential
equation on either side of $u=v$, imposing the two boundary
conditions, continuity, and the unit derivative jump. Its
eigenfunction expansion is uniformly absolutely convergent on
the closed square since its terms are bounded by
$(2/\alpha)/(x+(2j+1)^2)$. The spectral resolvent and the displayed
Green kernel agree in $L^2$ and hence by continuity everywhere.
At $u=v=\alpha$ this gives \eqref{parity:eq:S}; at $x=0$, continuity
gives $s(0)=\alpha$ and the odd reciprocal-square sum used above.

With Neumann conditions at both endpoints, the constant mode
contributes $1/(\alpha x)$ and the positive modes have
eigenvalues $x+4j^2$, $j\geq1$. The normalized endpoint values
have squares $2/\alpha$. The analogous Green kernel has endpoint
value $\coth(\alpha\sqrt x)/\sqrt x$. Multiplication by $x$ gives
the exact reciprocal series
\begin{equation}\label{parity:eq:reciprocal}
 \boxed{\frac1{s(x)}=\sqrt x\coth(\alpha\sqrt x)
  =\frac2\pi+\frac4\pi\sum_{j=1}^{\infty}\frac{x}{x+4j^2}.}
\end{equation}
The same uniformly convergent kernel argument justifies the
expansion for $x>0$. Both sides extend at zero to $2/\pi$.
Every summand is nonnegative. These identities are used on the
original axis $x>0$; no degree-dependent rescaling is made.

\section{Finite Stieltjes compression and native Gram approximation}

The following formulas apply to each original parity block with
$\nu=\nu_j$ and degree $n=n_j$. Put $v_n(x)=(1,x,\ldots,x^n)^T$,
\[
 H=\int v_nv_n^*\,d\nu,\quad
 X_H=\int x v_nv_n^*\,d\nu,\quad
 K=\int v_nv_n^*s\,d\nu,\quad
 K_x=\int x v_nv_n^*s\,d\nu.
\]
All four are finite positive definite matrices, and for the
zeroth parity block $X_H=H_{1,n}$ exactly. For $a>0$ define
$R(a)=\int v_nv_n^*/(x+a)\,d\nu$. The block Gram matrix of
$(x+a)^{1/2}v_n$ and $(x+a)^{-1/2}v_n$ is positive. Its Schur
complement gives
\begin{equation}\label{parity:eq:resolvent-compression}
 R(a)\succeq H(X_H+aH)^{-1}H.
\end{equation}
The positive series \eqref{parity:eq:S} therefore proves
\begin{equation}\label{parity:eq:finite-S}
 H^{1/2}s(H^{-1/2}X_HH^{-1/2})H^{1/2}
       \preceq K\preceq\alpha H.
\end{equation}
The matrix function in the lower bound is defined by the positive
spectral decomposition of the finite matrix. The summation is
justified because each inverse is at most $a^{-1}I$ and the
reciprocal-square series converges. This is a finite-moment
bound using the exact Christoffel moment matrix, not a replacement
of the full multiplication operator by its compression.

More directly useful for the unchanged original Gram is
\eqref{parity:eq:reciprocal}. Define the actual second-weight tilt
$d\widetilde\nu=s\,d\nu=r_jw_2\,dx$ and, for an integer $J\geq1$,
\begin{equation}\label{parity:eq:H-trunc}
 H^{[J]}=\frac2\pi K+\frac4\pi\sum_{j=1}^J
       \int\frac{x}{x+4j^2}v_nv_n^*\,d\widetilde\nu.
\end{equation}
Every summand is positive and $H^{[J]}$ is strictly positive
already from its first term. Integrating \eqref{parity:eq:reciprocal}
and using monotone convergence in every quadratic form gives
$H^{[J]}\uparrow H$. Since
$x/(x+4j^2)\leq x/(4j^2)$ and
$\sum_{j>J}j^{-2}\leq\int_J^\infty t^{-2}dt=1/J$, the full
matrix error bound is
\begin{equation}\label{parity:eq:H-error}
 \boxed{0\preceq H-H^{[J]}\preceq\frac1{\pi J}K_x
                                      \preceq\frac1{2J}X_H.}
\end{equation}
This is a certified approximation of the original native Gram,
not just a representation of an auxiliary tilted metric.

Only moments and negative-axis resolvent integrals are required
for a finite truncation. If
$\widetilde\mu_d=\int x^d\,d\widetilde\nu$ and
$\widetilde F(a)=\int(x+a)^{-1}\,d\widetilde\nu$, polynomial
division gives, for $d\geq1$,
\begin{equation}\label{parity:eq:resolvent-moments}
 \int\frac{x^d}{x+a}\,d\widetilde\nu
 =\sum_{r=0}^{d-1}(-a)^r\widetilde\mu_{d-1-r}
                  +(-a)^d\widetilde F(a).
\end{equation}
For $d=0$ the integral is $\widetilde F(a)$. This proves the
finite evaluation map, including every polynomial-division sign.
The positive integral form in \eqref{parity:eq:H-trunc} remains available
when cancellation in \eqref{parity:eq:resolvent-moments} would be
numerically inconvenient.

Apply these constructions to both blocks of
\eqref{parity:eq:parity-grams}. Denote their direct sums by $H^{[J]}$,
$K_x$ and $X_H$. The original source error is exactly the
congruence
\[
 0\preceq H_{S,N}^{\rm ar}-L_N^*H^{[J]}L_N
       \preceq\frac1{\pi J}L_N^*K_xL_N.
\]
For the unchanged onto observation $A$ in \eqref{parity:eq:observation},
the minimum-energy proof of \eqref{parity:eq:Q} shows that the map
$G\mapsto Q(G)=(AG^{-1}A^*)^{-1}$ is order-preserving on positive
source matrices: taking the minimum over the same set $Az=u$
preserves each quadratic-form inequality. Therefore
\begin{equation}\label{parity:eq:Q-bracket}
 \boxed{Q(H^{[J]})\preceq Q(H)\preceq
 Q\left(H^{[J]}+\frac1{\pi J}K_x\right).}
\end{equation}
Both endpoints converge to the original $Q(H)$: the finite
matrices converge to $H$, their positive inverses converge by
$B^{-1}-C^{-1}=B^{-1}(C-B)C^{-1}$, and the same holds after
the onto observation congruence. This proves a finite certified
source and observation calculation with the original classes and
Taylor unit untouched.

There is an entirely explicit degree-only bound for the error.
The original Gamma orthonormal recurrence has off-diagonal
$b_j=\sqrt{(j+1)(j+1/2)}\leq j+1$, so its two weighted shifts
give $\|yP\|_\sigma\leq2(N+1)\|P\|_\sigma$ on degree at most $N$.
Combining this with \eqref{parity:eq:CAI-finite} yields
\begin{equation}\label{parity:eq:L-bound}
 \|yP\|_{m_k}^2\leq\mathcal L_N\|P\|_{m_k}^2,
 \qquad \mathcal L_N=
       \frac{4A_k\Delta_N(N+1)^2}{a_k}.
\end{equation}
Under the parity map this is exactly $X_H\preceq\mathcal L_NH$.
Consequently for $J>\mathcal L_N/2$, with
$\varepsilon_J=\mathcal L_N/(2J)<1$, one has
\begin{equation}\label{parity:eq:relative}
 (1-\varepsilon_J)H\preceq H^{[J]}\preceq H,
 \qquad
 Q(H^{[J]})\preceq Q(H)
       \preceq(1-\varepsilon_J)^{-1}Q(H^{[J]}).
\end{equation}
This bound is generally coarse but is finite, positive and
explicit in the original CAI constants and degree. It does not
assume the desired arithmetic estimate.

\section{A second fully explicit bound using raw Pollaczek moments}

The literal ratio $r_0$ also admits a useful finite comparison
with the raw first literature weight. Put $L=M+1$ and
\begin{equation}\label{parity:eq:ratio-constants}
 c_0=\frac{a_kc_\Gamma\alpha}{2^{1/4}(1+2\alpha)},
 \qquad C_0=A_kC_\Gamma\alpha.
\end{equation}
Then for every $x>0$,
\begin{equation}\label{parity:eq:ratio-envelope}
 c_0(1+x)^{-L}\leq r_0(x)\leq C_0.
\end{equation}
To prove this, write $y=\sqrt x$ and retain the exact factor
\[
 e^{-\alpha y}\frac{\sinh(\alpha y)}y
    =\frac{1-e^{-2\alpha y}}{2y}.
\]
It is at most $\alpha$ and at least
$\alpha/(1+2\alpha y)$, since $1-e^{-u}\geq u/(1+u)$ for $u\geq0$.
The latter inequality follows from $e^u\geq1+u$.
Substituting the two CAI envelopes gives the upper bound and the
lower expression
\[
 \frac{a_kc_\Gamma\alpha(1+y^2)^{-M}}
                   {(1+2\alpha y)(1+y)^{1/2}}.
\]
Use $1+2\alpha y\leq(1+2\alpha)\sqrt{1+y^2}$ and
$(1+y)^{1/2}\leq2^{1/4}(1+y^2)^{1/4}$, followed by
$(1+y^2)^{-M-3/4}\geq(1+y^2)^{-M-1}$, to prove
\eqref{parity:eq:ratio-envelope} with the stated constants.

For parity $j=0,1$ use the raw measure
$d\tau^{(j)}=x^jw_1(x)dx$ and the degree $n=n_j$. Define the
explicit moment matrices
\[
 U_{j,n}=[\tau_{a+b+j}]_{a,b=0}^n,\qquad
 V_{j,n}=\left[\sum_{r=0}^L\binom Lr\tau_{a+b+j+r}\right]_{a,b=0}^n.
\]
All entries are the explicit values \eqref{parity:eq:raw-moments}.
The measures satisfy $d\nu_j=r_0\,d\tau^{(j)}$, so the same
block-Gram Schur-complement argument used in
\eqref{parity:eq:resolvent-compression}, now with weight $(1+x)^L$,
proves the actual native bound
\begin{equation}\label{parity:eq:raw-matrix-bound}
 \boxed{c_0 U_{j,n}V_{j,n}^{-1}U_{j,n}
             \preceq H_{j,n}\preceq C_0U_{j,n}.}
\end{equation}
Explicitly, the block matrix with blocks $V_{j,n},U_{j,n}$ and
$\int v_nv_n^*(1+x)^{-L}d\tau^{(j)}$ is a Gram matrix, giving
the lower rational-weight bound; \eqref{parity:eq:ratio-envelope}
then gives the displayed native inequality.
Taking block direct sums, followed by the unchanged congruence
$L_N$ or the unchanged observation $A$, transports these bounds
to the original source or to $Q(H)$, respectively. In particular,
if $G_-$ and $G_+$ are these two block bounds, then
\[
 Q(G_-)\preceq Q(H)\preceq Q(G_+).
\]
This gives a fully explicit finite comparison from raw literature
moments to the actual observation metric, with every density
ratio paid for. No raw multiple-orthogonal polynomial is being
used as a native orthogonal polynomial.

The same ratio envelope proves a typed map in the reverse
direction for the paper's actual multiple orthogonal polynomial
$P_{(n_1,n_2)}$. The function
\[
 f(x)=P_{(n_1,n_2)}(x)/r_0(x)
\]
belongs to $L^2(\nu_0)$, since its squared norm is
$\int |P|^2r_0^{-1}w_1dx$ and the inverse ratio is bounded by
$c_0^{-1}(1+x)^L$. Equations
\eqref{parity:eq:functional0}--\eqref{parity:eq:functionalS} give exactly
\[
 \int x^jf\,d\nu_0=0\ (0\leq j<n_1),\qquad
 \int x^jf\,s\,d\nu_0=0\ (0\leq j<n_2).
\]
This is an explicit receiver of the raw multiple-orthogonality
relations in the original even Hilbert sector. Its source image
is $r_0^{-1}\parityPol$, not the original polynomial source. For the
paper's diagonal rescaling
$Q_n(x)=(4n^2)^{-2n}P_{(n,n)}(4n^2x)$, the corresponding original
even-axis function is exactly
\[
 y\longmapsto\frac{(4n^2)^{2n}Q_n(y^2/(4n^2))}{r_0(y^2)}.
\]
This formula retains the complete coordinate change and multiplier;
no asymptotic statement about native $E_n$ or $O_n$ follows from
the paper's asymptotic theorem by omitting them.

\section{What the finite estimates do and do not complete}

The parity isometries, Christoffel maps and Gram congruences are
exact representations of the original objects. The nonzero
residual \eqref{parity:eq:residual-lower}, the rational truncation and
error bound \eqref{parity:eq:H-error}, the original observation brackets
\eqref{parity:eq:Q-bracket}--\eqref{parity:eq:relative}, and the explicit
raw-moment comparison \eqref{parity:eq:raw-matrix-bound} are additional
finite calculations and inequalities. Their constants and
receiving maps have all been evaluated above.

They preserve the original complex observed cross entries
\[
 B_i=(\Lambda a_i)^*Q(H)\Lambda v,
 \qquad
 \Phi_{B,N}=\frac{2\Re(\overline{B_1}B_2)}{\omega_{N,k}^{\rm ar}},
\]
with the actual boundary classes, terminal class, and physical
Taylor unit. The bounds certify source/observation approximation
and do not supply an evaluated sign, phase, or required
asymptotic of this last signed product. No such conclusion, and
no imported multiple-orthogonal asymptotic, is claimed.

\begin{thebibliography}{9}
\bibitem{parity:ALM} A. I. Aptekarev, G. L\'opez Lagomasino and
A. Mart\'{\i}nez--Finkelshtein,
\emph{Strong asymptotics for the Pollaczek multiple orthogonal
polynomials ensembles}, arXiv:1410.1261v1 (2014), Introduction:
the two unscaled weights and their positive Stieltjes ratio.
\url{https://arxiv.org/abs/1410.1261v1}.
The original author LaTeX is the source used; its multiple-orthogonal
asymptotic theorem is not assumed for the arithmetic weight.
\bibitem{parity:CAI} Split-Zero programme, \emph{The complete original
action relative to its arithmetic total}, CAI16--18 and its complete original
convolution providers. This is the programme calculation supplying
the actual arithmetic density comparison, not a human-literature
attribution for the Pollaczek weights.
\url{https://github.com/KokunoYumeto/zeta-function-research-reader/blob/2412f4aaa07c3fb19052cc0e537e93c2db9b0b1d/workbenches/splitzero-tandem/continuations/20260920-heat-literature-observed-class/providers/CAI_COMPLETE_PREREQUISITE.tex}.
\bibitem{parity:HG} Split-Zero programme, \emph{Heat growth and original
currents}, HG21 and its displayed proof, retaining every fixed
primary multiplicity and the original Gamma mass.
\url{https://github.com/KokunoYumeto/zeta-function-research-reader/blob/2412f4aaa07c3fb19052cc0e537e93c2db9b0b1d/workbenches/splitzero-tandem/continuations/20260920-heat-literature-observed-class/providers/HEAT_GROWTH_AND_ORIGINAL_CURRENTS.tex}.
The parity recurrences, Green-kernel expansions, finite Schur
inequalities and receiving maps are proved explicitly in this note;
no omitted asymptotic argument is attributed to these programme
providers.
\end{thebibliography}


\endgroup

\clearpage\begingroup
\part*{Growing-degree truncation and the full conductor}
\newcommand{\nativecutC}{\mathbb C}
\newcommand{\nativecutR}{\mathbb R}

This note carries the exact parity/Pollaczek receiver into the original
arithmetic multiplication estimate and the original conductor quotient.
The human input is the unscaled weight ratio of Aptekarev, L\'opez
Lagomasino and Mart\'{\i}nez--Finkelshtein \cite{nativecut:ALM}; all transport to
the arithmetic measure is proved in the complete companion \cite{nativecut:PP}.
The stronger truncation length below uses the already proved original
CAI21--25 estimate \cite{nativecut:CAI}, not a new comparison measure. No native
resolvent integral or signed phase is assigned an uncomputed value.

\section{Original objects and the proved finite truncation}
Retain the complete quartet divisor $h$, $g=2\xi$, the even source
$w_h(y)=|(g/h)(1/2+iy)|^2/(2\pi)$, and $m_k=w_h^{*k}$ with
$k=4\ell+1\ge9$, as in the companion. Their masses remain
$\mu_h^k$. For $S=c+iy$, $c=k/2$, the literal polynomial coordinate
map is $L_N=\Pi_NT_N$, where
$(T_N)_{rj}=\binom jr c^{j-r}i^r$ for $r\le j$ and $\Pi_N$ orders
the even powers first. The original source splits as
$f(y)=E(y^2)+yO(y^2)$ with the exact measures
\[
 d\nu_0(x)=m_k(\sqrt x)x^{-1/2}dx,\qquad d\nu_1=x\,d\nu_0.
\]
Let $H$ be their direct-sum monomial Gram at degrees
$\lfloor N/2\rfloor,\lfloor(N-1)/2\rfloor$. Then the original
$S$-coordinate Gram is $H_S=L_N^*HL_N$.

The literature multiplier and its exact reciprocal are
\[
 s(x)=\frac{\tanh(\pi\sqrt x/2)}{\sqrt x},\qquad
 \frac1{s(x)}=\frac2\pi+\frac4\pi\sum_{j\ge1}\frac{x}{x+4j^2}.
\]
Both identities, including convergence and constants, are proved from
the one-dimensional Green kernels in \cite{nativecut:PP}. For each parity block,
put $v(x)=(1,x,\ldots,x^n)^T$ and retain the actual tilt
$d\widetilde\nu=s\,d\nu$. Define
\[
 H^{[J]}=\frac2\pi\int vv^*d\widetilde\nu
 +\frac4\pi\sum_{j=1}^J\int\frac{x}{x+4j^2}vv^*d\widetilde\nu,
 \qquad X_H=\int xvv^*d\nu,
\]
and use direct sums for the two blocks. Positivity of every omitted
term and $\sum_{j>J}j^{-2}\le1/J$ give
\begin{equation}\tag{TR1}\label{nativecut:eq:tail}
 0\preceq H-H^{[J]}\preceq\frac1{\pi J}\int xvv^*s\,d\nu
                                  \preceq\frac1{2J}X_H.
\end{equation}
The last step uses $s(x)\le\pi/2$. All three integrals are actual
native integrals, not the raw Pollaczek moments.

\section{The original tensor bound sharpens the truncation length}
Here are the complete constants needed from CAI21--25 \cite{nativecut:CAI}:
\[
 \alpha=\pi/2,\ M=21+4m,\ \nu_*=\log(5/3),\ d_* =\alpha-\nu_*>0,
 \quad C_\sigma=\frac4{\sqrt{15}\sqrt{2\pi}},
\]
\[
 \begin{gathered}
 H_h=\frac{64C_\sigma(1+4M^2)^M}{d_*^2}
             \max_{b=3,4,5}\frac{M_\alpha^b}{a_b},\\
 T_h=\frac2{d_*}\{\log16+2M+3+\max(0,\log H_h)\},\qquad
 C_{\rm mult}(h)=\sqrt{T_h^2+1}.
 \end{gathered}
\]
The constants $a_b,M_\alpha$ are the original displayed constants
of the companion's CAI envelope, with the convolution order replaced
by $b$. CAI21 proves
\[
 \int_{|y|>R}y^2|P(y)|^2m_b(y)dy
 \le H_h(n+1)^{2M+2}16^ne^{-d_*R/2}\|P\|_{m_b}^2
 \quad(\deg P\le n,\ b=3,4,5).
\]
At $R=T_h(n+1)$ its coefficient is at most one, using
$\log(n+1)\le n$; inside the interval $y^2\le R^2$. Addition gives
$\|yP\|_{m_b}\le C_{\rm mult}(h)(n+1)\|P\|_{m_b}$.
Thus the tail result is used only at its established block orders.

Write $r_k=\lfloor k/3\rfloor$ and partition $k$ into $r_k$
summands $b_i\in\{3,4,5\}$, replacing one three by four or five
when necessary. Convolution proves the exact isometry
\[
 f(y)\longmapsto f(y_1+\cdots+y_{r_k})
 \quad\hbox{into }\bigotimes_iL^2(m_{b_i}(y_i)dy_i).
\]
Its image is the sum-dependent subspace and its mass is still
$\mu_h^k$. In each original orthonormal polynomial basis,
multiplication has zero diagonal by evenness and off-diagonal
coefficients $0<a_n^{(i)}\le C_{\rm mult}(h)n$: take the degree-$n$
coefficient in the preceding multiplication bound at degree $n-1$.

Compress the sum of multiplication operators to the product basis
of total degree at most $N+1$. At a vertex of degree at most $N$
its absolute row sum is at most $C_{\rm mult}(h)(2N+r_k)$; on the
outer face it has only downward entries and row sum at most
$C_{\rm mult}(h)(N+1)\le C_{\rm mult}(h)(2N+r_k)$. The symmetric
matrix norm is at most its maximal absolute row sum, by pairing
off-diagonal terms and using $2|uv|\le |u|^2+|v|^2$. Multiplication
of a degree-$N$ sum-dependent polynomial has degree at most $N+1$,
so no output is lost. Returning through the isometry proves
\begin{equation}\tag{TR2}\label{nativecut:eq:native}
 \|yP\|_{m_k}^2\le C_{\rm mult}(h)^2(2N+r_k)^2\|P\|_{m_k}^2.
\end{equation}
The companion also gives the full Gamma comparison
$\|yP\|_{m_k}^2\le[4A_k\Delta_N(N+1)^2/a_k]\|P\|_{m_k}^2$,
where $\Delta_N=[1+4(N+M)^2]^M$. Therefore retain both bounds as
\begin{equation}\tag{TR3}\label{nativecut:eq:eps}
 L_N^\sharp=\min\left\{\frac{4A_k\Delta_N(N+1)^2}{a_k},
                     C_{\rm mult}(h)^2(2N+r_k)^2\right\},
 \qquad \epsilon_{N,J}=\frac{L_N^\sharp}{2J}.
\end{equation}
Under the exact parity map, TR2 is $X_H\preceq L_N^\sharp H$.
For the actual integers $J>L_N^\sharp/2$, TR1 yields
\begin{equation}\tag{TR4}\label{nativecut:eq:relative}
 (1-\epsilon_{N,J})H\preceq H^{[J]}\preceq H.
\end{equation}
The sharper bound is quadratic in $N+k$ for fixed $h$. It does not
require an asymptotic estimate for the native orthogonal polynomials.

\section{Actual quotient metrics and signed endpoint errors}
Let $A$ be the original onto observation expressed in the parity
coordinates: $A=\Lambda J_{S,N}L_N^{-1}$. For a positive source
matrix $G$, the original attained minimum is
$Q_A(G)=(AG^{-1}A^*)^{-1}$. Its minimizing lift is
$G^{-1}A^*Q_A(G)u$; every other lift differs by a vector in $\ker A$
orthogonal to it for $G$. This proves the formula, and minimising
TR4 over the same fibre gives
\begin{equation}\tag{TR5}\label{nativecut:eq:Q}
 Q_{N,J}:=Q_A(H^{[J]})\preceq Q_N:=Q_A(H)
           \preceq(1-\epsilon_{N,J})^{-1}Q_{N,J}.
\end{equation}
Restriction to any specified original polynomial source subspace
preserves TR4 and the identical argument applies to its specified
onto observation. No different admitted relation space is substituted.

Write $d_N=\dim\operatorname{im}A$. The eigenvalues of
$Q_{N,J}^{-1/2}Q_NQ_{N,J}^{-1/2}$ lie between one and
$(1-\epsilon_{N,J})^{-1}$. Their product gives
\begin{equation}\tag{TR6}\label{nativecut:eq:det}
 0\le\log\frac{\det Q_N}{\det Q_{N,J}}
                 \le-d_N\log(1-\epsilon_{N,J}).
\end{equation}
This calculation proves an inequality for determinants in their
original frames; the square-root conjugation does not replace those
frames or norms. The empty determinant is one and satisfies TR6.

For the original endpoints $N_1,N_2,N_3,N_4$, retain
\[
 F=\log\frac{\det Q_{N_1}\det Q_{N_2}}
                  {\det Q_{N_3}\det Q_{N_4}},\qquad
 F_{\mathbf J}=\log\frac{\det Q_{N_1,J_1}\det Q_{N_2,J_2}}
                  {\det Q_{N_3,J_3}\det Q_{N_4,J_4}}.
\]
Put $e_i=-d_{N_i}\log(1-\epsilon_{N_i,J_i})$. Each individual
logarithmic error is nonnegative, so the complete signed enclosure is
\begin{equation}\tag{TR7}\label{nativecut:eq:four}
 -e_3-e_4\le F-F_{\mathbf J}\le e_1+e_2.
\end{equation}
For every specified tolerance $t>0$, choose actual integers
\begin{equation}\tag{TR8}\label{nativecut:eq:J}
 J_i\ge\left\lceil L_{N_i}^\sharp\max\{1,2d_{N_i}/t\}\right\rceil.
\end{equation}
Then $\epsilon_{N_i,J_i}\le1/2$ and
$-\log(1-\epsilon)\le2\epsilon$ give $e_i\le t/2$, hence
$|F-F_{\mathbf J}|\le t$. When $d_{N_i}=0$ its contribution is
zero and the same integer choice remains valid. At the packet
windows $N_i=O(q)$, $d_{N_i}\le q$ and fixed $h$, this uses
$O_h(q^3/t)$ resolvent terms for $0<t\le1$. Taking $t=1/k$ gives
a vanishing original four-endpoint truncation error. This is a
truncation-length bound, not a claim about numerical conditioning,
total complexity, or certification cost for the native integrals.
It does not evaluate the signed return or its limiting value.

\section{The same certified bound reaches the full conductor}
On the simple-quartet conductor domain, EQ1--21 \cite{nativecut:EQ} proves
$C=\bar C\Lambda$ and the full attained conductor metric
\[
 T_N=(\bar C Q_N^{-1}\bar C^*)^{-1},\qquad
 R_N=Q_N^{-1}\bar C^*T_N,
\]
with the same original source and with $\bar C$ onto. Apply the
same fixed-fibre minimum to TR5. Defining
$T_{N,J}=(\bar C Q_{N,J}^{-1}\bar C^*)^{-1}$ gives
\begin{equation}\tag{TR9}\label{nativecut:eq:conductor}
 T_{N,J}\preceq T_N\preceq
                    (1-\epsilon_{N,J})^{-1}T_{N,J}.
\end{equation}
The quotient identity also follows by direct substitution:
$T_{N,J}=(\bar CA(H^{[J]})^{-1}A^*\bar C^*)^{-1}$.
Thus the same truncation is applied to the original conductor
map, rather than to a separately chosen four-plane metric.
Restricting by its original cubic inclusion $I_W$ and coordinate
map $\Psi$ gives the identical relative bracket for
$G_N=\Psi^*I_W^*T_NI_W\Psi$. EQ16--19's exact pair-flip matrix
$B_J=O(I-D_J)O^{-1}E$ is independent of this metric. For every
original label vector $\alpha$ its actual defect energy obeys
\begin{equation}\tag{TR10}\label{nativecut:eq:energy}
 \alpha^*B_J^*G_{N,j}B_J\alpha
 \le\alpha^*B_J^*G_NB_J\alpha
 \le(1-\epsilon_{N,j})^{-1}\alpha^*B_J^*G_{N,j}B_J\alpha.
\end{equation}
Here the lower-case $j$ is the truncation length, distinct from
the two-element sign-flip set $J$; $G_{N,j}$ uses $T_{N,j}$.
The rank-two mixing map, all state labels, the Taylor unit and
the complete conductor kernel are unchanged. TR10 certifies its
positive defect energy, not the sign of a different projected
complex product. The terminal-class residual in EQ20--21 is
also retained; none of these bounds sets it to zero.

\begin{thebibliography}{9}
\bibitem{nativecut:ALM} A. I. Aptekarev, G. L\'opez Lagomasino and
A. Mart\'{\i}nez--Finkelshtein, \emph{Strong asymptotics for the
Pollaczek multiple orthogonal polynomials ensembles},
arXiv:1410.1261v1 (2014), Introduction, unscaled weights and
Stieltjes ratio. \url{https://arxiv.org/abs/1410.1261v1}.
The original author LaTeX was read in full by the integrating task.
\bibitem{nativecut:PP} Split-Zero programme, \emph{Exact native parity transport
into the Pollaczek pair and certified finite source--observation bounds},
complete companion \texttt{PARITY\_POLLACZEK\_NATIVE\_RECEIVER.tex},
equations \texttt{eq:reciprocal}, \texttt{eq:H-error} and
\texttt{eq:Q-bracket}, with their full proofs. This is the actual
arithmetic receiving calculation, not the author paper's asymptotics.
\bibitem{nativecut:CAI} Split-Zero programme, \emph{The complete original action
relative to its arithmetic total}, CAI21--25, in the complete provider
\texttt{CAI\_COMPLETE\_PREREQUISITE.tex}.
\url{https://github.com/KokunoYumeto/zeta-function-research-reader/blob/2412f4aaa07c3fb19052cc0e537e93c2db9b0b1d/workbenches/splitzero-tandem/continuations/20260920-heat-literature-observed-class/providers/CAI_COMPLETE_PREREQUISITE.tex}.
The original complete proof supplies the tail constant used above;
the total-degree transport and every new endpoint bound are proved here.
\bibitem{nativecut:EQ} Split-Zero programme, \emph{The eight signed states
through the original conductor minimum}, complete companion
\texttt{EIGHT\_STATE\_ORIGINAL\_MINIMUM.tex}, EQ1--21; independent
full derivation \texttt{EIGHT\_STATE\_QUOTIENT\_REDERIVATION.tex},
IQ1--29. Human provenance of the upstream factor-map construction
and exact source hashes are retained in those sources and the manifest.
\end{thebibliography}

\endgroup

\clearpage\begingroup
\part*{Eight signed states in the original observation minimum}
\newcommand{\eightminC}{\mathbb C}
\newcommand{\eightminPp}{\mathcal P}
\newcommand{\eightminim}{\operatorname{im}}
\newcommand{\eightminrank}{\operatorname{rank}}
\newcommand{\eightmindiag}{\operatorname{diag}}



This calculation receives the nonlinear ES--Fable construction through
the original finite conductor and the actual arithmetic observation
quotient. It evaluates the inherited eight-state minimum, its exact
correction to the four-dimensional source minimum, and the failure of
evaluation alone to retain the signed monodromy action. Every minimum
below is evaluated as a finite matrix of the unchanged original density.
No asymptotic allocation or arithmetic projected sign is assigned by
these finite identities. The human source of the underlying factor-map
construction is identified in Tao's exposition \cite{eightmin:Tao}; the marked
four-dimensional polynomial and its conductor receiver are the local
programme results \cite{eightmin:FC}, rather than results attributed to Tao.

\section{Original spaces and the actual factor through observation}
Use the simple-packet domain of OB1--3 \cite{eightmin:OB}:
\[
 \begin{gathered}
 \rho=\tfrac12+\delta+i\gamma,\quad0<\delta<\tfrac12,\quad\gamma>2,\\
 k\equiv1\pmod4,\ k\ge9,\quad j=k-8,\quad
 q=(k+1)^2,\quad q_-=(k-7)^2.
 \end{gathered}
\]
The fixed original period is in FC17's admitted five-orbit domain:
the five original transformed quadrics are pairwise nonassociate.
FC17 proves on this domain that the original symbol is nonzero and
its first nonzero moment has order at most eighty. This is the same
domain used by the receiving conductor; it is not an assertion about
an exceptional period outside that domain. All coefficients remain
unchanged. Write
\[
 \chi_k(S)=\prod_{a,b=0}^k
 [S-k/2-(2a-k)\delta-i(2b-k)\gamma],\qquad
 E_k=\eightminC[S]/(\chi_k),
\]
and define $\chi_j,E_j$ by the same expression with $k$ replaced by $j$.
The actual source density is $m_k(y)=w_h^{*k}(y)$ at $S=k/2+iy$,
where $h$ and $w_h$ retain OB1/LR2's complete original factors.
For every $N\ge q-1$, let $H_N$ be its positive Gram in the ordered
monomials $1,S,\ldots,S^N$. Positivity follows since a nonzero
polynomial has positive squared integral against this positive density.

Write $J_N:\eightminPp_N\to E_k$ for reduction modulo $\chi_k$ and
$\Lambda_k:E_k\to B_k$ for the original observation onto its image.
Set $\mathsf A_N=\Lambda_kJ_N$. OB3 proves the exact map
\[
 C:E_k\longrightarrow E_j,\quad C[P]=[T_AP],\qquad
 \ker\Lambda_k\subset\ker C.                                      \tag{EQ1}
\]
Thus the formula $\bar C(\Lambda_kx)=Cx$ is well-defined: two lifts
differ in $\ker\Lambda_k$, so have equal $C$-images. It is linear
and unique because $\Lambda_k$ is onto. In particular, the actual
polynomial map $\mathsf C_N=CJ_N$ satisfies
\[
 \bar C\mathsf A_N=\mathsf C_N.                                    \tag{EQ2}
\]
This constructs the needed receiving map. It does not assume that
$\Lambda_k$ descends through the different quotient
$\eightminPp_{v+3}/\eightminPp_{v-1}$ used in the four-plane construction.

Retain $v=\operatorname{ord}_0E_A\le80$ and the original moments
$\mu_n$. The finite-shift conductor has
\[
 T_A S^n=\sum_{r=0}^n\binom nr\mu_{n-r}(S')^r,
 \quad\mu_0=\cdots=\mu_{v-1}=0,\quad\mu_v\ne0.                    \tag{EQ3}
\]
The triangular map $\eightminPp_N\to\eightminPp_{N-v}$ is onto: its nonzero
diagonal from $S^{v+r}$ to $(S')^r$ is
$\binom{v+r}{r}\mu_v$. Moreover
$q-q_-=16k-48\ge96>80$, so $N-v\ge q_--1$.
Reduction onto $E_j$ is therefore onto. Consequently $\mathsf C_N$
and $\bar C$ are onto, with all of their original kernels retained.

The observation metric, the complete conductor quotient metric, and
their attained lifts are
\[
 \begin{split}
 Q_N&=(\mathsf A_NH_N^{-1}\mathsf A_N^*)^{-1},\\
 T_N&=(\mathsf C_NH_N^{-1}\mathsf C_N^*)^{-1}
       =(\bar C Q_N^{-1}\bar C^*)^{-1},\\
 L_N&=H_N^{-1}\mathsf C_N^*T_N:E_j\longrightarrow\eightminPp_N,\\
 R_N&=Q_N^{-1}\bar C^*T_N:E_j\longrightarrow B_k.
 \end{split}                                                       \tag{EQ4}
\]
All inverses exist by the surjectivity just proved. Direct substitution
gives $\mathsf C_NL_N=I$, $\bar CR_N=I$, and
$\mathsf A_NL_N=R_N$. If $x=L_Nu+d$ and $\mathsf C_Nd=0$, then
$d^*H_NL_Nu=d^*\mathsf C_N^*T_Nu=0$. Hence
\[
 \|x\|_{H_N}^2=u^*T_Nu+d^*H_Nd.                                  \tag{EQ5}
\]
This proves the full original minimum and uniqueness of its lift.
The identical calculation in $B_k$ proves
$R_N^*Q_NR_N=T_N$. Thus EQ4 is equality of successive actual
minima, not an asserted comparison of differently chosen norms.

The available Jacobi reconstruction \cite{eightmin:JC} evaluates every entry
used here from the original one-factor arithmetic transform. With its
unchanged matrices $D_N,C_N^{\rm Jac}$ one has
$H_N=D_N^*C_N^{\rm Jac}D_N$. Substitution in EQ4 keeps the physical
Taylor unit in $\mathsf A_N$ and retains the full kernel. No Gamma
replacement of $m_k$ is made in this substitution.

\section{Exact correction from the four-plane to the complete minimum}
Let $W=\eightminPp_3$ in the unchanged target variable $S'$ and
$I_W:W\to E_j$ be reduction modulo $\chi_j$. Since $q_-\ge4$,
this map is injective. Put
\[
 T_{W,N}=I_W^*T_NI_W.                                               \tag{EQ6}
\]
These are the norms of the complete original conductor lifts of
the actual cubic target values. The four-plane source minimum can
also be calculated exactly in the same arithmetic density.

Partition the principal Gram on $\eightminPp_{v+3}$ into the degrees below
$v$ and degrees $v,\ldots,v+3$:
\[
 H_{v+3}=\begin{pmatrix}H_{LL}&H_{LR}\\H_{RL}&H_{RR}\end{pmatrix},
 \quad S_v=H_{RR}-H_{RL}H_{LL}^{-1}H_{LR},
 \quad (B_v)_{rt}=\begin{cases}\binom{v+t}{r}\mu_{v+t-r},&r\le t,\\
 0,&r>t,\end{cases}                                                \tag{EQ7}
\]
for $0\le r,t\le3$. For $v=0$, the empty lower block is omitted
and $S_v=H_{RR}$. The diagonal of $B_v$ is
$\binom{v+t}{v}\mu_v$, so $B_v$ is invertible without altering
any moment. The unique minimum in $\eightminPp_{v+3}$ mapping to a cubic
with coefficient column $w$ is
\[
 L_{\rm low}w=\operatorname{incl}_{v+3,N}
 \begin{pmatrix}-H_{LL}^{-1}H_{LR}\\I_4\end{pmatrix}B_v^{-1}w,
 \qquad G_{\rm low}=B_v^{-*}S_vB_v^{-1}.                            \tag{EQ8}
\]
Indeed EQ3 forces the four upper coefficients to be $B_v^{-1}w$;
completing the square in the remaining lower coefficients proves
EQ8 and leaves all lower cross terms present. Define
\[
 D_N^{\rm lift}=L_{\rm low}-L_NI_W.
\]
Both lifts have the same $\mathsf C_N$-image, so
$\mathsf C_ND_N^{\rm lift}=0$. Applying EQ5 column by column gives
the exact positive correction
\[
 \boxed{G_{\rm low}-T_{W,N}
       =(D_N^{\rm lift})^*H_ND_N^{\rm lift}\succeq0.}                \tag{EQ9}
\]
In particular equality holds precisely when the two lift matrices
are identical. This answers the metric-identification question with
an explicit matrix of original moments and original maps. The four
source dimensions have not been used to remove the higher-degree
relations from the complete minimum.

\section{All eight states and their inherited observation minimum}
Use the actual four roots and their two factor signs in SE1--9
\cite{eightmin:FC}. Write an eight-state vector uniquely as
$\iota_+\alpha+\iota_-\beta$, retaining the factors $1/2$ in
$\iota_\pm e_r=(f_{r,+}\pm f_{r,-})/2$. The evaluated vector is
$z=E\alpha+O\beta$ in the original ES coordinates. SE7--8 proves
$O$ invertible and $\eightminrank E=3$. Retain the exact original maps
$Y:\eightminC^4\to W$, $\Psi=YK^{-1}:\eightminC^4\to W$ from FC32.
Thus the original target evaluation is $\Psi z$; no coordinate
column is silently declared orthonormal.

The actual lift into the original observation space is now defined,
rather than inferred from the existence of an abstract observation:
\[
 \mathcal O_N(\alpha,\beta)=R_NI_W\Psi(E\alpha+O\beta)\in B_k.
                                                                    \tag{EQ10}
\]
It satisfies $\bar C\mathcal O_N=I_W\Psi(E\alpha+O\beta)$.
Both $R_NI_W\Psi$ and $O$ are injective, so
\[
 \ker\mathcal O_N=\{(\alpha,-O^{-1}E\alpha):\alpha\in\eightminC^4\}.
                                                                    \tag{EQ11}
\]
This identifies exactly the four-dimensional graph kernel in the
original observation space, including the part invisible under sign
averaging.

Retain both labels and evaluation with the map
$\mathcal J(\alpha,\beta)=(Y\alpha,\Psi(E\alpha+O\beta))$.
The full original conductor minimum induces on its two target copies
the orthogonal sum $T_{W,N}\oplus T_{W,N}$. This is a specified
pullback from two copies of the actual minimum, not a previously
given probability measure on the eight abstract states. Define
\[
 H=Y^*T_{W,N}Y,\qquad G=\Psi^*T_{W,N}\Psi,
 \qquad
 \mathbb H_N=\begin{pmatrix}
 H+E^*GE&E^*GO\\O^*GE&O^*GO
 \end{pmatrix}.                                                     \tag{EQ12}
\]
The same proof applies to the original fixed Gamma target form
of FC56 by replacing $T_{W,N}$ with that actual form; the two
applications do not assert equality between those forms.
The map $(\alpha,\beta)\mapsto(\alpha,z)$ is invertible, with
$\beta=O^{-1}(z-E\alpha)$. Hence its full squared norm is exactly
\[
 \alpha^*H\alpha+z^*Gz.                                             \tag{EQ13}
\]
For the observation value $R_NI_W\Psi z_0$, its fibre is exactly
$z=z_0$. Minimisation in EQ13 therefore gives
\[
 \boxed{\min_{\mathcal O_N f=R_NI_W\Psi z_0}\|f\|_{\mathbb H_N}^2
       =z_0^*Gz_0=\|R_NI_W\Psi z_0\|_{Q_N}^2,\quad
       f_{\min}=\iota_-O^{-1}z_0.}                                  \tag{EQ14}
\]
This evaluates the inherited quotient with no additional norm loss
from the four graph-kernel directions. Its relation to the restricted
source cost is exactly EQ9, not an assumption that that correction
vanishes. Indeed $G$ equals
$\Psi^*G_{\rm low}\Psi-\Psi^*(D_N^{\rm lift})^*H_ND_N^{\rm lift}\Psi$.

For completeness the entire inverse of EQ12 is
\[
 \mathbb H_N^{-1}=
 \begin{pmatrix}
 H^{-1}&-H^{-1}E^*O^{-*}\\
 -O^{-1}EH^{-1}&O^{-1}(G^{-1}+EH^{-1}E^*)O^{-*}
 \end{pmatrix}.                                                     \tag{EQ15}
\]
Multiplication by EQ12 verifies this formula; alternatively it
follows by inverting the two triangular maps in EQ13.
In particular $[E\ O]\mathbb H_N^{-1}[E\ O]^*=G^{-1}$.
For any further actual linear observation $A z$ onto its image,
the inherited minimum is $(AG^{-1}A^*)^{-1}$, with unique minimizing
state $(\alpha,\beta)=(0,O^{-1}G^{-1}A^*(AG^{-1}A^*)^{-1}b)$.
This last identity is a fully evaluated finite extension; it does not
postulate an unidentified arithmetic observation map.

\section{The nonzero action defect on those exact fibres}
SM1--14 \cite{eightmin:FC} proves that the actual monodromy contains every
even change of the four factor signs. For a two-element subset
$J\subset\{1,2,3,4\}$ write $D_J$ for its diagonal sign change.
In the coordinates $(\alpha,z)$ of EQ13 its action is precisely
\[
 (\alpha,z)\longmapsto(\alpha,B_J\alpha+L_Jz),\quad
 L_J=OD_JO^{-1},\qquad B_J=O(I-D_J)O^{-1}E.                           \tag{EQ16}
\]
This is the restriction of the complete action SE16, derived directly
by substituting $z=E\alpha+O\beta$ and $\beta\mapsto D_J\beta$.
The canonical section of evaluation is $s(z)=(0,O^{-1}z)$.
On that section $g_Js(z)=s(L_Jz)$. On a general state in the
same evaluation fibre, the extra term is exactly $B_J\alpha$.
The action is therefore not determined by the original observed
value alone whenever this term survives the observation.

Here its size is also an original metric expression. Put
$f_\alpha=(\alpha,-O^{-1}E\alpha)\in\ker\mathcal O_N$. Then
\[
 \mathcal O_N(g_J f_\alpha)=R_NI_W\Psi B_J\alpha,\quad
 \|f_\alpha\|_{\mathbb H_N}^2=\alpha^*H\alpha,\quad
 \|\mathcal O_N(g_J f_\alpha)\|_{Q_N}^2
       =\alpha^*B_J^*GB_J\alpha.                                   \tag{EQ17}
\]
Thus its exact squared norm on the forgotten kernel is
$\lambda_{\max}(H^{-1/2}B_J^*GB_JH^{-1/2})$. There is no replacement
of this expression by the determinant of the nonlinear map.

Every pair flip has a rank-two defect in EQ17. Indeed SE1 gives
the first row of $O$ as $(a_1,a_2,a_3,a_4)$ with all $a_r\ne0$;
SE8 gives $\eightminim E=\{x:x_1=0\}$. Therefore
\[
 \eightminim(O^{-1}E)=\{x:\textstyle\sum_r a_rx_r=0\}.                     \tag{EQ18}
\]
Projection of this hyperplane onto either prescribed pair of
coordinates is onto: choose an omitted coordinate, whose coefficient
is nonzero, to solve the one equation. Consequently
\[
 \eightminim B_J=O\operatorname{span}\{e_r:r\in J\},\quad\eightminrank B_J=2,
 \qquad\sum_{|J|=2}\eightminim B_J=\eightminC^4.                                  \tag{EQ19}
\]
All equalities use the literal original $a_r$, without fixing their
signs. Since $R_NI_W\Psi$ is injective, each of the six pair flips
has rank-two failure to preserve $\ker\mathcal O_N$.
Thus none induces an action on its evaluation quotient.

More generally, for every nonzero linear observation $A$ on the
four evaluated coordinates, at least one actual pair flip has
$AB_J\ne0$. Otherwise EQ19 would imply $A=0$. Choose $\alpha$
with $AB_J\alpha\ne0$. The states $0$ and $f_\alpha$ have the same
original observation zero, but their transformed observations differ
by $AB_J\alpha$. This proves non-descent even for set maps, not only
linear candidates. The full exact receiver is EQ16 on both retained
coordinates, and its missing term after forgetting the labels is
the explicitly evaluated EQ17. Restriction to the chosen section
is a different, stated operation: $L_{J_1}L_{J_2}=OD_{J_1}D_{J_2}O^{-1}$
does give a representation of the sign subgroup there, but it does
not repair the representative dependence just exhibited.

\section{Scope of the receiving result}
The nonlinear example is therefore carried through an actual map
into $B_k$, with its entire inherited quotient metric calculated.
The finite observation values supplied here form precisely the
four-dimensional image $R_NI_W W$. No additional arithmetic
eigenclass has been inserted into it. In particular the lower
terminal cardinal $v_-$ in OB3 has degree $q_--1$: for $k=9$ this
is three, while for $k\ge13$ it is at least 35. Reduction modulo
$\chi_j$ has a unique representative of degree below $q_-$, so
in the latter range $v_-$ does not belong to $I_WW$. The exact
bridge to every original target class nevertheless remains EQ4;
for the terminal class it gives the attained lift
$R_N(a_{88}v_-)$ of the literal conductor image
$Cv_+=a_{88}v_-$. This statement retains the original scalar and
does not assert equality with the original class $\Lambda_kv_+$.
Their difference lies in $\ker\bar C$, and its exact norm identity is
\[
 \|\Lambda_kv_+\|_{Q_N}^2
 =|a_{88}|^2v_-^*T_Nv_-
   +\|\Lambda_kv_+-R_N(a_{88}v_-)\|_{Q_N}^2.                        \tag{EQ20}
\]
The cross term vanishes by the same orthogonality as EQ5, now in
$B_k$. This also gives the actual attained original polynomial.
Define $M_N=H_N^{-1}\mathsf A_N^*Q_N$, so
$\mathsf A_NM_N=I$, $M_N^*H_NM_N=Q_N$, and EQ2--4 gives
$M_NR_N=L_N$. Therefore
\[
 M_N\Lambda_kv_+
 =L_N(a_{88}v_-)+M_N(\Lambda_kv_+-R_N(a_{88}v_-)),\qquad
 \mathsf C_NM_N(\Lambda_kv_+-R_N(a_{88}v_-))=0.                       \tag{EQ21}
\]
Its two summands are $H_N$-orthogonal by EQ5 and have precisely
the two energies in EQ20. Thus the result supplies both the exact
full-class return and the complete residual, without mistaking the
eight-state four-plane for the entire growing packet.

\begin{thebibliography}{9}
\bibitem{eightmin:Tao} Terence Tao, \emph{A digestion of the Jacobian conjecture
counterexample}, 21 July 2026.
\url{https://terrytao.wordpress.com/2026/07/21/a-digestion-of-the-jacobian-conjecture-counterexample/}.
Human exposition of the Alp\"oge--Fable construction; it does not
contain the programme's marked four-dimensional receiver or EQ1--20.
\bibitem{eightmin:FC} The Clankers, ES--Fable marked polynomial construction,
canonical ES reader version 69, theorem
\texttt{explicit-four-time-Jacobian-collision}; programme continuation,
\emph{Fable to original conductor}, FC26--71, GF1--28, SE1--18,
SM1--19. The full exact source versions and reading coverage are in
the accompanying source receipt. Published predecessor:
\url{https://github.com/KokunoYumeto/zeta-function-research-reader/tree/369f6ea9e0312bdf348ca9efd294735efc99c51d}.
\bibitem{eightmin:OB} Programme original observed-class return, OB1--3;
original WCF finite-shift conductor and AC coefficient providers.
\url{https://github.com/KokunoYumeto/zeta-function-research-reader/tree/2412f4aaa07c3fb19052cc0e537e93c2db9b0b1d/workbenches/splitzero-tandem/continuations/20260920-heat-literature-observed-class}.
\bibitem{eightmin:JC} Programme, \emph{The actual arithmetic multiplier in
Gamma Jacobi coordinates}, 20 September 2026, equations
\texttt{eq:sourcegram}, \texttt{eq:quotient}; applying
Franti\v{s}ek \v{S}tampach and Pavel \v{S}\v{t}ov\'{i}\v{c}ek,
\emph{Spectral representation of some weighted Hankel matrices and
orthogonal polynomials from the Askey scheme}, arXiv:1811.05820v1.
The finite minimum proofs EQ4--20 are given here in full.
\end{thebibliography}

\endgroup

\clearpage\begingroup
\part*{Independent full minimum and monodromy-defect proof}
\newcommand{\indminC}{\mathbb C}
\newcommand{\indminim}{\operatorname{im}}
\newcommand{\indmindiag}{\operatorname{diag}}
\newcommand{\indminrank}{\operatorname{rank}}
\newcommand{\indminSpan}{\operatorname{span}}
\newcommand{\indminid}{\operatorname{id}}



\section{Exact scope and the retained source maps}
This note derives finite-dimensional identities from the actual eight-state
maps SE1--18 in \texttt{FABLE\_TO\_ORIGINAL\_CONDUCTOR.tex}, read together
with the explicit monodromy generators SM1--19 of that source.  All quotient
minima below retain the supplied positive Hermitian forms.  No inner product
on the eight abstract labels is assigned independently of the stated original
two-copy map.  In particular, the factors $1/2$ in the following embeddings
are retained:
\[
 \iota_+e_j=(f_{j,+}+f_{j,-})/2,\qquad
 \iota_-e_j=(f_{j,+}-f_{j,-})/2.                       \tag{IQ1}
\]
The actual matrices $E=(e_1\ \cdots\ e_4)$ and $O=(o_1\ \cdots\ o_4)$
have $\indminrank E=3$, $\det O\ne0$, and
\[
 F=O^{-1}E,\qquad
 \indminim F=\{x\in\indminC^4:a_1x_1+\cdots+a_4x_4=0\},\quad a_j\ne0. \tag{IQ2}
\]
These are exactly SE7, SE8 and SE18, with the original branches $a_j$.
The four-plane target maps $Y,\Psi:\indminC^4\longrightarrow W$ are invertible.
For $h=\iota_+\alpha+\iota_-\beta$ write
$\mathcal E h=E\alpha+O\beta$ and set
\[
 z=E\alpha+O\beta,\qquad
 \mathcal J_W h=(Y\alpha,\Psi z),\qquad
 H=Y^*\langle\ ,\ \rangle_WY,\quad
 G=\Psi^*\langle\ ,\ \rangle_W\Psi.                  \tag{IQ3}
\]
Here the expressions for $H,G$ denote the Gram matrices in the displayed
coordinate bases.  Thus $H,G$ are positive definite, with every original
cross entry and mass present.  The exact inverse and inherited norm are
\[
 h=\iota_+\alpha+\iota_-O^{-1}(z-E\alpha),\qquad
 \|h\|_{\mathcal J_W}^2=\alpha^*H\alpha+z^*Gz.         \tag{IQ4}
\]

Let $\Lambda:W\to Z$ be a specified linear observation, with $Z=\indminim\Lambda$;
choose any fixed coordinate basis of this image.  Put $A=\Lambda\Psi$.
Then $A$ is an $m\times4$ matrix of rank $m=\dim Z$, and the eight-state
observation is
\[
 \mathcal R h=\Lambda\Psi(E\alpha+O\beta)=Az.          \tag{IQ5}
\]
This is a finite observation calculation, not an assertion that a different
programme observation descends through a polynomial quotient.  In particular,
an observation defined on a larger original polynomial space cannot be applied
to $\mathcal P_{v+3}/\mathcal P_{v-1}$ without proving that it annihilates
$\mathcal P_{v-1}$.  IQ5 uses only the expressly specified map on $W$.  A
conductor receiver from a larger source must be constructed on its original
observation image before using IQ5.

\section{Evaluation of the full inherited minimum}
For $m>0$ define
\[
 Q=(AG^{-1}A^*)^{-1},\qquad R_G=G^{-1}A^*Q,\qquad
 P_G=R_GA.                                          \tag{IQ6}
\]
The inverse exists: if $b\ne0$, surjectivity of $A$ makes $A^*b\ne0$,
so $b^*AG^{-1}A^*b>0$.  Direct multiplication gives
\[
 AR_G=I_Z,\quad P_G^2=P_G,\quad P_G^*G=GP_G,
 \quad\ker P_G=\ker A.                              \tag{IQ7}
\]
Every vector with $Az=b$ has a unique expression $z=R_Gb+n$, $n\in\ker A$.
Its cross term vanishes because
\[
 (R_Gb)^*Gn=b^*QA n=0,\qquad
 (R_Gb)^*G(R_Gb)=b^*Qb.
\]
Consequently the complete original minimum is evaluated, not merely bounded:
\[
 \begin{split}
 \|h\|_{\mathcal J_W}^2
   &=b^*Qb+\alpha^*H\alpha+n^*Gn,\\
 \min_{\mathcal Rh=b}\|h\|_{\mathcal J_W}^2
   &=b^*Qb
    =\min_{\Lambda w=b}\|w\|_W^2,\\
 s_W(b)&=\iota_-O^{-1}G^{-1}A^*Qb.
 \end{split}                                       \tag{IQ8}
\]
The minimum is unique: both omitted positive terms vanish exactly when
$\alpha=0$ and $n=0$.  The eight-state minimizing vector has coefficients
$\beta_j/2,-\beta_j/2$ at $f_{j,+},f_{j,-}$, respectively; IQ1 has not been
replaced by a Euclidean normalization.  For $m=0$ the only observation value
is zero, and its unique minimum is $h=0$.

For clarity, the Gram matrix in the original $(\alpha,\beta)$ coordinates is
\[
 M=\begin{pmatrix}
 H+E^*GE&E^*GO\\ O^*GE&O^*GO
 \end{pmatrix},\qquad
 M^{-1}=\begin{pmatrix}
 H^{-1}&-H^{-1}E^*O^{-*}\\
 -O^{-1}EH^{-1}&O^{-1}(G^{-1}+EH^{-1}E^*)O^{-*}
 \end{pmatrix}.                                    \tag{IQ9}
\]
These identities follow by multiplying the triangular coordinate change
$S=\left(\begin{smallmatrix}I&0\\E&O\end{smallmatrix}\right)$, its inverse
$S^{-1}=\left(\begin{smallmatrix}I&0\\-O^{-1}E&O^{-1}\end{smallmatrix}\right)$,
and $M=S^*\indmindiag(H,G)S$.  In particular the full evaluation row
$\mathsf R=A(E\ O)$ satisfies
\[
 \mathsf R M^{-1}\mathsf R^*=AG^{-1}A^*.              \tag{IQ10}
\]
Thus the cancellation of the even-label block follows from the actual
attained minimum; it is not a deletion of that block before minimization.
The complete observation kernel is
\[
 \ker\mathcal R=
 \{\iota_+\alpha+\iota_-O^{-1}(n-E\alpha):
             \alpha\in\indminC^4,\ n\in\ker A\},          \tag{IQ11}
\]
of dimension $8-m$.  It retains the four-dimensional evaluation graph
kernel ($n=0$) and all additional unseen target directions.

\section{The exact source analogue and conductor receiver}
Let $L=T_A|_U:U\to W$ be the original invertible four-plane conductor,
and retain the original maps $LX=Y$, $L\Phi=\Psi$.  Put
\[
 H_U=X^*\langle\ ,\ \rangle_UX,\qquad
 G_U=\Phi^*\langle\ ,\ \rangle_U\Phi,
 \qquad\mathcal J_Uh=(X\alpha,\Phi z).               \tag{IQ12}
\]
The source observation used here is expressly $\Lambda L:U\to Z$.
Its coordinate matrix is $\Lambda L\Phi=A$; it is not an unidentified
quotient of an ambient observation.  Repeating the proof IQ6--8 gives
\[
 Q_U=(AG_U^{-1}A^*)^{-1},\quad R_U=G_U^{-1}A^*Q_U,
 \quad s_U(b)=\iota_-O^{-1}R_Ub,\quad
 \min_{\Lambda Lu=b}\|u\|_U^2=b^*Q_Ub.              \tag{IQ13}
\]
The conductor $L\oplus L$ intertwines the two total maps and induces the
identity map on the common observation value $b\in Z$.  This identity has
source metric $Q_U$ and target metric $Q$; no equality of those metrics is
assumed.  Its two original minimizing lifts differ by the explicit vector
\[
 d(b)=L\Phi R_Ub-\Psi R_Gb
     =\Psi(R_U-R_G)b\in\ker\Lambda,\qquad
 \|L\Phi R_Ub\|_W^2=b^*Qb+\|d(b)\|_W^2.            \tag{IQ14}
\]
Indeed $A(R_U-R_G)=0$, and the orthogonality calculation in IQ8 applies.
If $\sigma_{\min},\sigma_{\max}$ are the actual extremal singular values
of $L$ for these original norms, then minimization of
$\sigma_{\min}^2\|u\|_U^2\le\|Lu\|_W^2
\le\sigma_{\max}^2\|u\|_U^2$ over the same fibre gives
\[
 \sigma_{\min}^2 Q_U\preceq Q\preceq\sigma_{\max}^2 Q_U. \tag{IQ15}
\]
This proves the exact metric receiver and its inherited bounds without
asserting a new estimate for the original conductor singular values.

\section{Full monodromy action and every observation defect}
For a signed root permutation, retain the matrices
$\Pi e_j=e_{\pi(j)}$ and $Re_j=\sigma_j e_{\pi(j)}$.
The source calculation SE16 gives
\[
 B_g=ORO^{-1},\quad C_g=E\Pi-B_gE,\qquad
 g:(\alpha,z)\longmapsto(\Pi\alpha,B_gz+C_g\alpha).
                                                               \tag{IQ16}
\]
Substitution of $z=E\alpha+O\beta$ proves this formula directly from
$(\alpha,\beta)\mapsto(\Pi\alpha,R\beta)$.  In particular it retains
the complete mixing block, and has the exact group law.  For
$z=R_Gb+n$, $n\in\ker A$, its observed value is
\[
 \mathcal R(gh)=T_gb+AB_gn+AC_g\alpha,\qquad
 T_g=AB_gR_G.                                         \tag{IQ17}
\]
This is the entire defect from an operation on the quotient observation.
It proves the exact criterion: $g$ induces a linear map on that quotient
if and only if
\[
 AB_g(\ker A)=0\quad\hbox{and}\quad AC_g=0.          \tag{IQ18}
\]
Necessity follows by applying IQ17 to the two independently variable
parts of IQ11; sufficiency follows by the same formula.  In that case
the induced map is exactly $T_g$, independently of the chosen representative.

Every two-sign flip actually occurs in the monodromy calculated in
SM10--14.  Indeed the square of the displayed adjacent four-cycle flips
the two adjacent signs.  The three vectors $(1,1,0,0),(0,1,1,0),(0,0,1,1)$
span all even sign vectors over $\mathbb F_2$, proving the assertion for
every pair, without treating an arbitrary signed permutation as a proved
geometric loop.

For the flip $D_J$ of a two-element set $J$, IQ16 becomes
\[
 B_J=OD_JO^{-1},\qquad
 C_J=O(I-D_J)F,
 \qquad \indminim C_J=O\indminSpan\{e_j:j\in J\}.              \tag{IQ19}
\]
To prove the last equality, projection of $\indminim F$ onto the coordinates in
$J$ is onto: prescribe those two coordinates, set one of the remaining
coordinates to zero, and solve for the other using its nonzero $a_k$ in
IQ2.  Since $I-D_J$ multiplies exactly the coordinates in $J$ by two,
the assertion follows with the scalar two retained.  Therefore
\[
 \indminim(AC_J)=\indminSpan\{AOe_j:j\in J\},\quad
 \indminrank(AC_J)=\dim\indminSpan\{AOe_j:j\in J\},\quad
 \sum_{|J|=2}\indminim(AC_J)=Z.                            \tag{IQ20}
\]
The final equality holds since $O$ is invertible and $A$ is onto $Z$.
In particular, if the observation is nonzero, some $AOe_j$ is nonzero,
and every pair containing that index gives $AC_J\ne0$.  Thus no nonzero
linear observation of evaluation alone carries the full even-sign
monodromy action.  This statement is stronger than a failure of one
chosen four-dimensional presentation: it is proved for every nonzero
linear receiving observation of that evaluation.

Here is an explicit zero-class witness, including the branch factors.
Choose $j\in J$ with $AOe_j\ne0$ and an index $k\notin J$.  Put
\[
 x=e_j-(a_j/a_k)e_k\in\indminim F,\qquad
 F\alpha=x,\qquad
 h_0=\iota_+\alpha-\iota_-F\alpha.                 \tag{IQ21}
\]
Existence of $\alpha$ follows from the proved image identity IQ2.
Then $z(h_0)=0$, but
\[
 \mathcal Rh_0=0,\qquad
 \mathcal R(g_Jh_0)=2AOe_j\ne0.                    \tag{IQ22}
\]
There is no unidentified cancellation or asymptotic term in this witness.

\section{The defect measured in the inherited quotient metric}
Restrict the full defect to the evaluation graph kernel in IQ11, namely
$h_0(\alpha)=\iota_+\alpha-\iota_-F\alpha$.  Its original norm is
$\|h_0(\alpha)\|^2=\alpha^*H\alpha$.  Its image under an even flip
has observed quotient norm exactly
\[
 \|\mathcal R(g_Jh_0(\alpha))\|_Q^2
   =\alpha^*C_J^*A^*QA C_J\alpha.                  \tag{IQ23}
\]
Hence the complete squared operator norm of this mixing defect, from
the inherited graph-kernel norm to the inherited quotient norm, is
\[
 \|\mathcal Rg_J|_{\ker\mathcal E}\|^2
 =\lambda_{\max}
   \bigl(H^{-1/2}C_J^*A^*QA C_JH^{-1/2}\bigr).       \tag{IQ24}
\]
The formula follows from the Rayleigh quotient after the invertible
substitution $v=H^{1/2}\alpha$.  Positivity of $Q$ proves that IQ24
is strictly positive exactly when $AC_J\ne0$.  In particular, IQ21
gives the explicit lower bound
\[
 \|\mathcal Rg_J|_{\ker\mathcal E}\|^2
 \ge \frac{4\,(AOe_j)^*Q(AOe_j)}{\alpha^*H\alpha}>0. \tag{IQ25}
\]
If the smallest possible denominator is desired, let
$V=\indminim F$ carry the restriction of the original coordinate Euclidean
form solely for defining the adjoint in this explicit solution, and set
$B=(FH^{-1}F^*)|_V:V\to V$.  This operator is positive definite on
$V$: for $v\in V$, $v^*Bv=\|H^{-1/2}F^*v\|^2$, and
$F^*v=0$ with $v\in\indminim F$ forces $v=0$.  Then
\[
 \alpha_*=H^{-1}F^*B^{-1}x,\qquad
 F\alpha_*=x,\qquad
 \alpha_*^*H\alpha_*=x^*B^{-1}x.                   \tag{IQ26}
\]
For any other preimage $\alpha_*+t$, $Ft=0$ and
$\alpha_*^*Ht=x^*B^{-1}Ft=0$; this proves the claimed minimizing
denominator.  The resulting original-metric lower bound in IQ25 is
therefore $4\|AOe_j\|_Q^2/(x^*B^{-1}x)$.  The auxiliary coordinate
adjoint in IQ26 changes no norm in IQ23--25.

The canonical minimizing section in IQ8 also has an exact action defect.
Since its $\alpha$ coordinate is zero, the difference between its
monodromy image and the minimizing lift of the new observation is
\[
 g s_W(b)-s_W(T_gb)
 =\iota_-O^{-1}(I-P_G)B_gR_Gb\in\ker\mathcal R.     \tag{IQ27}
\]
The same orthogonality as in IQ8 evaluates its full cost:
\[
 \|g s_W(b)\|_{\mathcal J_W}^2
 =b^*T_g^*QT_gb+
  \bigl((I-P_G)B_gR_Gb\bigr)^*G
       \bigl((I-P_G)B_gR_Gb\bigr).                 \tag{IQ28}
\]
These formulas do not imply that the group preserves the original
positive forms.  They compute its actual failure to preserve the
minimizing section, with all original cross entries retained.
Finally, since $B_gB_h=B_{gh}$, the matrices of these compressed
operations have the exactly evaluated composition defect
\[
 T_gT_h-T_{gh}
   =-AB_g(I-P_G)B_hR_G.                            \tag{IQ29}
\]
This expression can vanish for particular observations; for example a
four-dimensional invertible observation has $P_G=I$.  Even in that
case the independent mixing defect IQ23 can be nonzero, and IQ22
still proves that these compressed operations do not define the
action of the entire eight-state space modulo evaluation.  Equations
IQ17, IQ23 and IQ27 specify all three exact receiving maps required
to compare these presentations.

\section{Provenance and limits of the derivation}
The input source is the programme proof
\texttt{work/rh\_counterfactual\_20260913/fable\_conductor\_20260919/}
\texttt{FABLE\_TO\_ORIGINAL\_CONDUCTOR.tex}, specifically SE1--18 and
SM1--19.  The finite-dimensional proofs IQ1--29 above are supplied in
full rather than attributed to a human source not used here.  They
prove the exact attained metric transfer and the complete observation
defect.  They do not assert that the ambient arithmetic observation
descends through an unverified source quotient, that the nonlinear
cover changes an original conductor kernel, or that the positive
defect yields an RH sign estimate.

\endgroup

% END COMPLETE PARITY AND CONDUCTOR RECEIVERS
\end{document}
